%% 4-6 坂道と三角比（第13回：組 V4-4。ページ制限なし。偶数ページ）
%% ほかの巻から指されている形：
%%   5-17「4-6 の『斜面の重力を2つに分ける』」→ 道具⑥（mg sinθ と mg cosθ。なぜ斜面の角と同じ θ が出てくるかを 4-1 の角で）
%%   3-back のつぎの巻へ「坂の上の荷物にかかる力を、2つに分ける」（はてな？）
\makeatletter
% 式を、まん中に1行で置く  \bkivgline{…}
\newcommand\bkivgline[1]{\par\addvspace{3pt}{\centering #1\par}\addvspace{3pt}\noindent\ignorespaces}
% 説明の段の見出し（丸数字＋太字の1行）  \bkivgstep{番号}{見出し}
\newcommand\bkivgstep[2]{\par\addvspace{6pt}\Needspace{4\baselineskip}\noindent
  {\bkcircled[bkOrange]{#1}\enspace\sffamily\bfseries\color{bkOrange!75!black}#2}\par\nobreak\vspace{2pt}}
% 小見出し（それだけで1行）  \bkivghead{見出し}
\newcommand\bkivghead[1]{\par\addvspace{4pt}\Needspace{5\baselineskip}\noindent\textbf{#1}\par\nobreak}
% 荷物（箱）。下の辺のまん中を (x,y) に置き、角 #2 だけ回す  \bkivgbox{(x,y)}{角}{色}
\newcommand\bkivgbox[3]{\begin{scope}[shift={#1}, rotate=#2]
  \filldraw[fill=#3!35, draw=#3!70!black, line width=0.7pt, rounded corners=0.6mm] (-6,0) rectangle (6,9);
  \draw[#3!70!black, line width=0.5pt] (-6,6.2) -- (6,6.2);\end{scope}}
% 台車（荷物をのせた箱に車輪）。下の辺のまん中を (x,y) に置き、角 #2 だけ回す  \bkivgcart{(x,y)}{角}{色}
\newcommand\bkivgcart[3]{\begin{scope}[shift={#1}, rotate=#2]
  \filldraw[fill=bkGray!60, draw=bkInk, line width=0.5pt] (-3.6,1.4) circle (1.4) (3.6,1.4) circle (1.4);
  \filldraw[fill=#3!35, draw=#3!70!black, line width=0.7pt, rounded corners=0.6mm] (-6,2.8) rectangle (6,10.8);
  \draw[#3!70!black, line width=0.5pt] (-6,7.6) -- (6,7.6);\end{scope}}
% 直角のしるし  \bkivgright{角の点}{1本めの向き(度)}{大きさ mm}
\newcommand\bkivgright[3]{\draw[bkInk, line width=0.5pt] ($#1+(#2:#3)$) -- ++(#2+90:#3) -- ++(#2+180:#3);}
\makeatother

\bkchapter{4}{6}{坂道と三角比}{\bklevelrange{2}{3}}{坂道と三角比}
\bkmichishirube{3-5,4-1,4-3,4-4}{4-7,4-8,5-17}

\begin{goal}
\begin{itemize}[label={\color{\bkTT}●}, leftmargin=1.6\zw, itemsep=0pt, topsep=1pt]
\item \bklevel{2}\enspace 直角三角形で、どの角から見るかに気をつけて、$\sin$・$\cos$・$\tan$ を求められる。
\item \bklevel{3}\enspace 三角比と電卓を使って、辺の長さや角度を求められる。坂の傾きを角度に直せる。
\item \bklevel{3}\enspace 斜面にのった物の重力を、$mg\sin\theta$ と $mg\cos\theta$ の2つに分けられる。
\end{itemize}
\end{goal}

\begin{hatena}
\noindent\begin{minipage}[c]{0.54\linewidth}
荷物をのせた台車を、坂の途中で止めておきたい。平らな所なら、手をはなしても台車は動きません。でも坂の上では、手をはなすと、台車は坂を下っていきます。

台車を引っぱる重力は、まっすぐ下向きです。それなのに、なぜ坂にそって動くのでしょう。台車を止めておくには、坂にそって、どれだけの力で支えればよいのでしょう。
\end{minipage}\hfill
\begin{minipage}[c]{0.43\linewidth}
\centering
\begin{tikzpicture}[x=1mm, y=1mm]
  \def\t{28}
  \fill[bkGray!25] (0,0) -- (52,0) -- (52,{52*tan(\t)}) -- cycle;
  \draw[bkInk!80, line width=0.8pt] (0,0) -- (52,0) -- (52,{52*tan(\t)}) -- cycle;
  \coordinate (S) at ({32*cos(\t)},{32*sin(\t)});
  \bkivgcart{(S)}{\t}{bkOrange}
  \coordinate (G) at ($(S)+({\t+90}:6.8)$);
  \fill[bkInk] (G) circle (0.8);
  \draw[bk arrow, bkBlue] (G) -- ++(0,-16);
  \node[font=\sffamily\footnotesize, text=bkBlue!80!black, anchor=west] at ($(G)+(1.2,-11)$) {重力};
  \draw[bk arrow, bkRed, dashed, line width=1pt] (G) -- ++({\t+180}:10);
  \node[font=\sffamily\footnotesize\bfseries, text=bkRed!80!black, anchor=south] at ($(G)+({\t+180}:10)+(0,1.6)$) {？};
  \draw[bkInk, line width=0.5pt] (8,0) arc (0:\t:8);
  \node[font=\sffamily\footnotesize, anchor=north] at (6,-1.5) {坂の角};
\end{tikzpicture}
\end{minipage}
\end{hatena}

\noindent 坂のような「ななめ」の形は、くらしの中にたくさんあります。車いすのスロープ、山の道路、屋根、すべり台。この章では、ななめの形を、直角三角形の辺の比で表します。その比を使うと、坂の角度から高さを計算したり、重力を2つに分けたりできます。

%% ---------------- やってみよう・思い出そう ----------------
\begin{yattemiyou}[同じ角の直角三角形を3つかこう]
\noindent\begin{minipage}[c]{0.56\linewidth}
\textbf{用意するもの}\par
ノート、定規、分度器、電卓
\begin{enumerate}[label={\bkcircled[bkGreen]{\arabic*}}, leftmargin=2.4\zw, itemsep=1pt, topsep=3pt]
\item 横の辺が\qty{4}{cm}で、左はしの角が\qty{30}{\degree}の直角三角形をかこう（右下を直角にする）。
\item 横の辺を\qty{6}{cm}、\qty{8}{cm}にして、同じ\qty{30}{\degree}の直角三角形をかこう。
\item たての辺と、ななめの辺の長さをはかり、「たて÷ななめ」「横÷ななめ」「たて÷横」を電卓で計算しよう。
\item 角を\qty{40}{\degree}にして、もう1つかいてみよう。
\end{enumerate}
\end{minipage}\hfill
\begin{minipage}[c]{0.41\linewidth}\centering
\begin{tikzpicture}[x=1mm, y=1mm]
  \foreach \b/\c in {40/bkBlue, 30/bkGreen, 20/bkOrange} {
    \draw[\c!80!black, line width=0.8pt] (0,0) -- (\b,0) -- (\b,{\b*tan(30)}) -- cycle;}
  \draw[bkInk, line width=0.5pt] (8,0) arc (0:30:8);
  \node[font=\sffamily\footnotesize] at (15:11.5) {\ang{30}};
  \foreach \b in {20,30,40} {\draw[bkInk, line width=0.4pt] (\b-2.2,0) -- (\b-2.2,2.2) -- (\b,2.2);}
  \node[font=\sffamily\footnotesize, anchor=north] at (20,-1.5) {4};
  \node[font=\sffamily\footnotesize, anchor=north] at (30,-1.5) {6};
  \node[font=\sffamily\footnotesize, anchor=north] at (40,-1.5) {8\,cm};
\end{tikzpicture}
\end{minipage}
\par\vspace{4pt}
ソラさんがはかった結果です（長さは\qty{0.1}{cm}まで、比は小数第2位まで）。
\begin{bkcenter}
\begin{bktable}{colspec={Q[c,m,wd=11mm] Q[c,m,wd=11mm] Q[c,m,wd=11mm] Q[c,m,wd=13mm] Q[c,m,wd=21mm] Q[c,m,wd=21mm] Q[c,m,wd=17mm]}, colsep=2pt, rowsep=1.2pt, row{1}={font=\sffamily\footnotesize}}
角 & 横 & たて & ななめ & たて÷ななめ & 横÷ななめ & たて÷横 \\
\ang{30} & 4.0 & 2.3 & 4.6 & 0.50 & 0.87 & 0.58 \\
\ang{30} & 6.0 & 3.5 & 6.9 & 0.51 & 0.87 & 0.58 \\
\ang{30} & 8.0 & 4.6 & 9.2 & 0.50 & 0.87 & 0.58 \\
\ang{40} & 6.0 & 5.0 & 7.8 & 0.64 & 0.77 & 0.83 \\
\end{bktable}
\end{bkcenter}
\end{yattemiyou}

\noindent \ang{30}の3つの三角形は、大きさがちがうのに、3つの比がほとんど同じです（0.51のような小さなずれは、はかるときのずれです）。角を\ang{40}にすると、比が変わりました。\textbf{直角三角形の辺の比は、大きさではなく、角だけで決まる}ようです。この章では、この比に名前をつけて使います。

\begin{omoidasou}
\bkrecall{3-5}{傾き＝上がる高さ÷横に進む長さ（たての変化量÷横の変化量）。道路ではこう配ともいう。}
\bkrecall{4-1}{三角形の3つの角（内角）をたすと\ang{180}。平行な2本の直線に1本の直線が交わってできる同位角・錯角は等しい。}
\bkrecall{4-3}{2つの角がそれぞれ等しい三角形は相似。相似な図形では、対応する辺の比が等しい。直角三角形の辺の比は、角だけで決まる。}
\bkrecall{4-4}{直角をはさむ辺が $a$、$b$、斜辺が $c$ なら $a^2+b^2=c^2$（三平方の定理）。\\三角定規の辺の比は $1:1:\sqrt2$ と $1:2:\sqrt3$。}
\bkrecall{4-2}{電卓は DEG（度）と RAD（ラジアン）を切りかえて使う。表計算の \texttt{SIN} はラジアンで計算する。\\極座標 $(r,\ \theta)$：中心からの距離 $r$ と、$x$ 軸の正の向きから左回りの角 $\theta$ で場所を言う。}
\end{omoidasou}

%% ---------------- 数学の道具 ----------------
\bksection{doguu}{数学の道具}

\bksubsection{1}{坂の急さの3つの言い方}

\leavevmode\label{bk:4-6:kyuu}町の中で、坂の急さは、いろいろな言い方で書かれています。

\begin{bkfigure}
\begin{tikzpicture}[x=1mm, y=1mm]
  % スロープ（1/12）
  \fill[bkGray!20] (0,0) -- (36,0) -- (36,3) -- cycle;
  \draw[bkInk!80, line width=0.7pt] (0,0) -- (36,0) -- (36,3) -- cycle;
  \node[font=\sffamily\footnotesize, anchor=north, align=center] at (18,-2) {車いすのスロープ\\「$\dfrac{1}{12}$ より急にしない」};
  % 道路の標識 10%
  \begin{scope}[shift={(66,6)}]
  \fill[bkYellow, rounded corners=1mm, rotate around={45:(0,0)}] (-9,-9) rectangle (9,9);
  \draw[bkInk, line width=0.8pt, rounded corners=1mm, rotate around={45:(0,0)}] (-8,-8) rectangle (8,8);
  \fill[bkInk] (-6.5,-3) -- (6.5,-3) -- (6.5,3) -- cycle;
  \node[font=\sffamily\bfseries\footnotesize, text=bkInk] at (-1.5,4.5) {10\%};
  \end{scope}
  \node[font=\sffamily\footnotesize, anchor=north, align=center] at (66,-8) {道路の標識\\「こう配10\%」};
  % スキー場 30°
  \fill[bkBlue!10] (96,0) -- (126,0) -- (96,{30*tan(30)}) -- cycle;
  \draw[bkInk!80, line width=0.7pt] (96,0) -- (126,0) -- (96,{30*tan(30)}) -- cycle;
  \draw[bkInk, line width=0.5pt] (118,0) arc (180:150:8);
  \node[font=\sffamily\footnotesize, anchor=south east] at (114,0.6) {\ang{30}};
  \node[font=\sffamily\footnotesize, anchor=north, align=center] at (111,-2) {スキー場のコース\\「最大斜度\ang{30}」};
\end{tikzpicture}
\bkfigcaption{同じ「坂の急さ」を、分数・百分率（\%）・角度の3つの言い方で表している。}
\end{bkfigure}

\bkivgstep{1}{傾き（3-5）}
上がる高さ÷横に進む長さです。スロープの「$\dfrac{1}{12}$」は、横に\qty{12}{m}進むと\qty{1}{m}上がる、という意味です。傾きは $1\div12=0.083\cdots$ です。

\bkivgstep{2}{こう配（\%）}
道路の標識は、傾きを百分率（\%）で書いたものです。傾きを100倍すると\%になります（2-5）。「こう配10\%」は、横に\qty{100}{m}進むと\qty{10}{m}上がる坂で、傾きは $10\div100=0.1$ です。スロープの $\dfrac{1}{12}$ は、$0.083\times100$ で約8.3\%です。

\bkivgstep{3}{角度}
坂と水平な線のあいだの角を、分度器ではかります。スキー場の「斜度\ang{30}」は、この角が\ang{30}という意味です。

\par\addvspace{4pt}
3つの言い方を、表にまとめます。角度は、図を大きくかいて分度器ではかった値です（あとで計算でも求めます）。
\begin{bkcenter}
\begin{bktable}{colspec={Q[l,m,wd=34mm] Q[c,m,wd=24mm] Q[c,m,wd=24mm] Q[c,m,wd=24mm]}, rowsep=1.5pt, row{1}={font=\sffamily\small}}
 & 傾き & こう配 & 角度 \\
スロープ $\dfrac{1}{12}$ & 0.083 & 8.3\% & 約\ang{4.8} \\
道路の標識 & 0.1 & 10\% & 約\ang{5.7} \\
横とたてが同じ坂 & 1 & 100\% & \ang{45} \\
横1でたて2の坂 & 2 & 200\% & 約\ang{63.4} \\
\end{bktable}
\bktabcaption{傾きとこう配は、100倍するだけ。角度は、かんたんには計算できない。}
\end{bkcenter}

\noindent\begin{minipage}[c]{0.54\linewidth}
表の下の2行を見ましょう。傾きが1から2へ、2倍になりました。角度は\ang{45}の2倍の\ang{90}になるでしょうか。\par\vspace{2pt}
なりません。\ang{63.4}です。傾き2の坂は「横に1進むと2上がる」坂で、まだかべ（\ang{90}）ではないからです。
\end{minipage}\hfill
\begin{minipage}[c]{0.42\linewidth}\centering
\begin{tikzpicture}[x=0.93mm, y=0.93mm]
  \draw[bkGray, line width=0.6pt, -{Stealth[length=4pt]}] (-2,0) -- (44,0);
  \draw[bkGray, line width=0.6pt, -{Stealth[length=4pt]}] (0,-2) -- (0,44);
  \draw[bkBlue, line width=1.2pt] (0,0) -- (38,19);
  \draw[bkGreen!70!black, line width=1.2pt] (0,0) -- (30,30);
  \draw[bkOrange, line width=1.2pt] (0,0) -- (19,38);
  \node[font=\sffamily\footnotesize, text=bkBlue!80!black, anchor=west] at (38.8,19) {傾き0.5};
  \node[font=\sffamily\footnotesize, text=bkGreen!50!black, anchor=west] at (30.8,30) {傾き1};
  \node[font=\sffamily\footnotesize, text=bkOrange!80!black, anchor=west] at (19.8,38) {傾き2（\ang{63.4}）};
  \draw[bkBlue!80!black, line width=0.5pt] (13,0) arc (0:26.57:13);
  \node[font=\sffamily\footnotesize, text=bkBlue!80!black] at (13.3:17.6) {\ang{26.6}};
  \draw[bkGreen!50!black, line width=0.5pt] (22,0) arc (0:45:22);
  \node[font=\sffamily\footnotesize, text=bkGreen!50!black] at (36:26) {\ang{45}};
  \draw[bkOrange!80!black, line width=0.5pt] (7,0) arc (0:63.43:7);

\end{tikzpicture}
\end{minipage}
\par\vspace{4pt}
傾きをどんなに大きくしても、角度は\ang{90}に近づくだけです。\textbf{傾きと角度は、比例しません}。傾きと角度を行き来するには、新しい道具が必要です。それが、この章で作る\textbf{三角比}です。

\bksubsection{2}{直角三角形の辺の比：$\sin$・$\cos$・$\tan$}

\leavevmode\label{bk:4-6:teigi}\bkyattemiyouref では、角が\ang{30}なら、三角形の大きさを変えても、辺の比が同じになりました。なぜでしょう。2通りで考えます。

\bkivgstep{1}{相似で考える（4-3）}
\noindent\begin{minipage}[c]{0.5\linewidth}
角が $\theta$（シータ。4-2）の直角三角形を、大きさを変えて3つかき、$\theta$ の所を重ねます。3つの三角形は、どれも角が $\theta$ と\ang{90}です。三角形の角をたすと\ang{180}なので（4-1）、のこりの角は、どれも $\ang{180}-\ang{90}-\theta=\ang{90}-\theta$ です。\par\vspace{2pt}
3つの角がそれぞれ等しいので、3つの三角形は相似です（4-3）。相似な図形は、対応する辺の比が等しいので、「たて÷ななめ」などの比は、3つとも同じです。
\end{minipage}\hfill
\begin{minipage}[c]{0.46\linewidth}\centering
\begin{tikzpicture}[x=1mm, y=1mm]
  \def\t{32}
  \foreach \b/\c in {46/bkBlue, 33/bkGreen, 20/bkOrange} {
    \fill[\c!10] (0,0) -- (\b,0) -- (\b,{\b*tan(\t)}) -- cycle;}
  \foreach \b/\c in {46/bkBlue, 33/bkGreen, 20/bkOrange} {
    \draw[\c!80!black, line width=0.9pt] (0,0) -- (\b,0) -- (\b,{\b*tan(\t)}) -- cycle;
    \draw[bkInk, line width=0.4pt] (\b-2.2,0) -- (\b-2.2,2.2) -- (\b,2.2);}
  \draw[bkInk, line width=0.5pt] (8,0) arc (0:\t:8);
  \node[font=\sffamily\footnotesize] at ({\t/2}:11) {$\theta$};
  \foreach \b in {20,33,46} {\draw[bkInk, line width=0.5pt] ($(\b,{\b*tan(\t)})+(180+\t:5)$) arc (180+\t:270:5);}
  \node[font=\sffamily\footnotesize, anchor=west] at (47.5,{46*tan(\t)-4}) {$\ang{90}-\theta$};
\end{tikzpicture}
\end{minipage}

\bkivgstep{2}{坂をのぼって考える}
\noindent\begin{minipage}[c]{0.5\linewidth}\centering
\begin{tikzpicture}[x=1mm, y=1mm]
  \def\t{30}
  \fill[bkGray!15] (0,0) -- ({60*cos(\t)},0) -- ({60*cos(\t)},{60*sin(\t)}) -- cycle;
  \draw[bkInk!80, line width=0.9pt] (0,0) -- ({60*cos(\t)},{60*sin(\t)});
  \draw[bkInk!80, line width=0.7pt] (0,0) -- ({60*cos(\t)},0);
  \foreach \s/\lab in {20/1, 40/2, 60/3} {
    \fill[bkOrange] ({\s*cos(\t)},{\s*sin(\t)}) circle (1);
    \draw[bkBlue, line width=1pt] ({\s*cos(\t)},0) -- ({\s*cos(\t)},{\s*sin(\t)});
    \node[font=\sffamily\footnotesize, anchor=south east, inner sep=1pt] at ({\s*cos(\t)-0.8},{\s*sin(\t)+0.4}) {\lab\,m};}
  \node[font=\sffamily\footnotesize, text=bkBlue!80!black, anchor=west] at ({20*cos(\t)+0.8},{10*sin(\t)-1.6}) {0.5};
  \node[font=\sffamily\footnotesize, text=bkBlue!80!black, anchor=west] at ({40*cos(\t)+0.8},{20*sin(\t)-1}) {1};
  \node[font=\sffamily\footnotesize, text=bkBlue!80!black, anchor=west] at ({60*cos(\t)+0.8},{30*sin(\t)}) {1.5};
  \draw[bkInk, line width=0.5pt] (7,0) arc (0:\t:7);
  \node[font=\sffamily\footnotesize] at ({\t/2}:10) {\ang{30}};
\end{tikzpicture}
\end{minipage}\hfill
\begin{minipage}[c]{0.46\linewidth}
\ang{30}の坂を、坂にそって歩きます。\qty{1}{m}歩くと高さは\qty{0.5}{m}、\qty{2}{m}歩くと\qty{1}{m}、\qty{3}{m}歩くと\qty{1.5}{m}になります。同じ坂をのぼりつづけるので、歩いた長さと上がった高さは比例します（2-15）。だから「高さ÷歩いた長さ」は、どこではかっても0.5です。
\end{minipage}

\par\vspace{4pt}
どちらの考え方でも、\textbf{直角三角形の辺の比は、角 $\theta$ だけで決まる}ことがわかりました。そこで、この比に名前をつけます。

\bkivghead{辺の名前：どの角から見るか}
\noindent\begin{minipage}[c]{0.52\linewidth}
直角三角形の3つの辺に、角 $\theta$ から見た名前をつけます。$\theta$ の所に目を置いて、三角形をながめてください。
\begin{itemize}[leftmargin=1.2\zw, itemsep=2pt, topsep=3pt, label={\color{\bkTT}●}]
\item \textbf{斜辺}（しゃへん）：直角の向かいにある辺。いちばん長い辺。
\item \textbf{対辺}（たいへん）：$\theta$ の向かいにある辺。目の正面に見える辺。
\item \textbf{隣辺}（りんぺん）：$\theta$ のとなりの辺のうち、斜辺でないほう。
\end{itemize}
\end{minipage}\hfill
\begin{minipage}[c]{0.45\linewidth}\centering
\begin{tikzpicture}[x=1mm, y=1mm]
  \coordinate (A) at (0,0); \coordinate (B) at (40,0); \coordinate (C) at (40,25);
  \fill[bkTheme!12] (A) -- (B) -- (C) -- cycle;
  \draw[bkRed, line width=1.6pt] (B) -- (C);
  \draw[bkBlue, line width=1.6pt] (A) -- (C);
  \draw[bkGreen!70!black, line width=1.6pt] (A) -- (B);
  \bkivgright{(B)}{90}{2.6}
  \draw[bkInk, line width=0.5pt] (8,0) arc (0:32:8);
  \node[font=\sffamily\footnotesize] at (16:11) {$\theta$};
  \node[font=\sffamily\footnotesize, text=bkBlue!80!black, anchor=south east] at (19,13.5) {斜辺};
  \node[font=\sffamily\footnotesize, text=bkRed!80!black, anchor=west] at (41.5,12.5) {対辺};
  \node[font=\sffamily\footnotesize, text=bkGreen!50!black, anchor=north] at (20,-1.5) {隣辺};
  % 目
  \fill[white] (-6,1) ellipse (3 and 1.8);
  \draw[bkInk, line width=0.5pt] (-6,1) ellipse (3 and 1.8);
  \fill[bkInk] (-5.2,1) circle (0.9);
\end{tikzpicture}
\end{minipage}

\par\vspace{4pt}
3つの辺の比のうち、よく使う3つに名前をつけます。
\begin{bkcenter}
\begin{bktable}{colspec={Q[c,m,wd=24mm] Q[c,m,wd=30mm] Q[c,m,wd=30mm] Q[l,m,wd=36mm]}, rowsep=2.5pt, row{1}={font=\sffamily\small}}
記号 & 読み方 & 式 & 坂なら \\
$\sin\theta$ & サイン & $\dfrac{\text{対辺}}{\text{斜辺}}$ & 高さ÷ななめの長さ \\
$\cos\theta$ & コサイン & $\dfrac{\text{隣辺}}{\text{斜辺}}$ & 横÷ななめの長さ \\
$\tan\theta$ & タンジェント & $\dfrac{\text{対辺}}{\text{隣辺}}$ & 高さ÷横（傾き） \\
\end{bktable}
\bktabcaption{三角比の決め方（定義。4-1）。どれも $\theta$ から見た辺の名前で決める。}
\end{bkcenter}

\noindent $\sin\theta$ を\textbf{正弦}（せいげん）、$\cos\theta$ を\textbf{余弦}（よげん）、$\tan\theta$ を\textbf{正接}（せいせつ）ともいいます。3つをまとめて\textbf{三角比}（さんかくひ）といいます。\bkidx{さんかくひ}{三角比}\bkidx{さいん}{$\sin$（サイン）}$\sin\theta$ は「$\sin$ と $\theta$ のかけ算」ではありません。「角 $\theta$ のサイン」という1つの数です。$f(x)$（3-3）のように、$\theta$ を入れると比が1つ出てくる、と考えるとよいでしょう。

\bkyattemiyouref の表から、\ang{30}と\ang{40}の三角比は、次のようになります。
\bkivgline{\begin{tabular}{@{}l@{\qquad}l@{\qquad}l@{}}$\sin\ang{30}=0.50$ & $\cos\ang{30}=0.87$ & $\tan\ang{30}=0.58$\\[2pt] $\sin\ang{40}=0.64$ & $\cos\ang{40}=0.77$ & $\tan\ang{40}=0.83$\end{tabular}}

\begin{seikaku}
\textbf{三角比は、ここでは $\ang{0}<\theta<\ang{90}$ の角で決めた}\par
直角三角形の直角でない角は、\ang{0}より大きく\ang{90}より小さい角（鋭角（えいかく））です。だから、ここで決めた三角比は、鋭角のときだけの決め方です。\ang{90}より大きい角の三角比は、⑧で決め直します。
\end{seikaku}

\bkivghead{斜辺を1にすると、たては $\sin\theta$、横は $\cos\theta$}
\noindent\begin{minipage}[c]{0.52\linewidth}
長さ\qty{1}{m}の棒を、地面から角 $\theta$ だけ立てかけます。棒の先の高さは、対辺÷斜辺＝$\sin\theta$ で、斜辺が1なので、高さはそのまま $\sin\theta$（m）です。棒の先までの横の長さは $\cos\theta$（m）です。\par\vspace{2pt}
棒の長さが\qty{5}{m}なら、三角形は相似なまま5倍になるので、高さは $5\sin\theta$、横は $5\cos\theta$ です。
\end{minipage}\hfill
\begin{minipage}[c]{0.45\linewidth}\centering
\begin{tikzpicture}[x=1mm, y=1mm]
  \def\t{35}\def\L{44}
  \fill[bkTheme!12] (0,0) -- ({\L*cos(\t)},0) -- ({\L*cos(\t)},{\L*sin(\t)}) -- cycle;
  \draw[bkInk, line width=1.6pt] (0,0) -- ({\L*cos(\t)},{\L*sin(\t)});
  \draw[bkOrange, line width=1.4pt] ({\L*cos(\t)},0) -- ({\L*cos(\t)},{\L*sin(\t)});
  \draw[bkBlue, line width=1.4pt] (0,0) -- ({\L*cos(\t)},0);
  \bkivgright{({\L*cos(\t)},0)}{90}{2.4}
  \draw[bkInk, line width=0.5pt] (7,0) arc (0:\t:7);
  \node[font=\sffamily\footnotesize] at ({\t/2}:10) {$\theta$};
  \node[font=\sffamily\footnotesize, anchor=south east] at ({\L*cos(\t)/2-1},{\L*sin(\t)/2+1}) {1};
  \node[font=\sffamily\footnotesize, text=bkOrange!80!black, anchor=west] at ({\L*cos(\t)+1.5},{\L*sin(\t)/2}) {$\sin\theta$};
  \node[font=\sffamily\footnotesize, text=bkBlue!80!black, anchor=north] at ({\L*cos(\t)/2},-1.5) {$\cos\theta$};
\end{tikzpicture}
\end{minipage}

\par\vspace{4pt}
同じことを、式でも確かめます。$\sin\theta=\dfrac{\text{対辺}}{\text{斜辺}}$ の両辺に斜辺をかけると（2-14）、
\bkivgline{\fbox{対辺＝斜辺×$\sin\theta$}\qquad \fbox{隣辺＝斜辺×$\cos\theta$}\qquad \fbox{対辺＝隣辺×$\tan\theta$}}
となります。3つめは、$\tan\theta=\dfrac{\text{対辺}}{\text{隣辺}}$ の両辺に隣辺をかけたものです。この「かけ算の形」は、このあと何度も使います。

\bkivghead{$\sin\theta$・$\cos\theta$ は、極座標を $(x,\ y)$ に直す道具}\label{bk:4-6:kyoku}
4-2の極座標を思い出しましょう。点Pを、原点Oからの距離 $r$ と、$x$ 軸の正の向きから左回りの角 $\theta$ で、$(r,\ \theta)$ と言いました。同じ点Pを、3-2の座標 $(x,\ y)$（右へいくつ、上へいくつ）で言うと、どうなるでしょう。

\noindent\begin{minipage}[c]{0.5\linewidth}
Pから $x$ 軸へ垂直な線を下ろすと、直角三角形ができます。斜辺はOPで長さ $r$、角 $\theta$ から見た隣辺が $x$、対辺が $y$ です。だから、上の「かけ算の形」から、次の式ができます。
\end{minipage}\hfill
\begin{minipage}[c]{0.46\linewidth}\centering
\begin{tikzpicture}[x=1mm, y=1mm]
  \def\t{40}\def\R{40}
  \draw[bkGray, line width=0.6pt, -{Stealth[length=4pt]}] (-4,0) -- (40,0) node[anchor=west, font=\sffamily\footnotesize, text=bkInk] {$x$};
  \draw[bkGray, line width=0.6pt, -{Stealth[length=4pt]}] (0,-4) -- (0,32) node[anchor=south, font=\sffamily\footnotesize, text=bkInk] {$y$};
  \fill[bkTheme!12] (0,0) -- ({\R*cos(\t)},0) -- ({\R*cos(\t)},{\R*sin(\t)}) -- cycle;
  \draw[bkInk, line width=1.3pt] (0,0) -- ({\R*cos(\t)},{\R*sin(\t)});
  \draw[bkBlue, line width=1.4pt] (0,0) -- ({\R*cos(\t)},0);
  \draw[bkOrange, line width=1.4pt] ({\R*cos(\t)},0) -- ({\R*cos(\t)},{\R*sin(\t)});
  \bkivgright{({\R*cos(\t)},0)}{90}{2.4}
  \fill[bkInk] ({\R*cos(\t)},{\R*sin(\t)}) circle (1);
  \node[font=\sffamily\footnotesize, anchor=south west, inner sep=1pt] at ({\R*cos(\t)+0.6},{\R*sin(\t)+0.4}) {P};
  \draw[bkInk, line width=0.5pt] (7,0) arc (0:\t:7);
  \node[font=\sffamily\footnotesize] at ({\t/2}:10.5) {$\theta$};
  \node[font=\sffamily\footnotesize, anchor=south east] at ({\R*cos(\t)/2-1},{\R*sin(\t)/2+1}) {$r$};
  \node[font=\sffamily\footnotesize, text=bkBlue!80!black, anchor=north] at ({\R*cos(\t)/2},-1.5) {$x=r\cos\theta$};
  \node[font=\sffamily\footnotesize, text=bkOrange!80!black, anchor=west] at ({\R*cos(\t)+1.6},{\R*sin(\t)/2}) {$y=r\sin\theta$};
  \node[font=\sffamily\footnotesize, anchor=north east, inner sep=1pt] at (-0.6,-0.6) {O};
\end{tikzpicture}
\end{minipage}
\bkivgline{\fbox{$x=r\cos\theta$\qquad $y=r\sin\theta$}}
いまは、$\ang{0}<\theta<\ang{90}$ の向きで考えています（⑧で\ang{180}まで、4-8でどの向きにも広げます）。極座標 $(r,\ \theta)$ がわかれば、$r$ に $\cos\theta$ と $\sin\theta$ をかけるだけで、$(x,\ y)$ が求められます。\textbf{$\sin\theta$ と $\cos\theta$ は、距離と向きの言い方 $(r,\ \theta)$ と、横とたての言い方 $(x,\ y)$ をつなぐ橋}なのです。

数でも確かめます。4-2のロボットのうでは、肩から\qty{30}{cm}、向き\ang{40}の所に手がありました（極座標 $(30,\ \ang{40})$）。\bkyattemiyouref ではかった $\cos\ang{40}=0.77$、$\sin\ang{40}=0.64$ を使うと、
\bkivgline{$x=30\times0.77=23.1$\qquad $y=30\times0.64=19.2$}
手は、肩から右へ約\qty{23}{cm}、上へ約\qty{19}{cm}の所です。三平方の定理（4-4）で確かめると、$23.1^2+19.2^2=533.61+368.64=902.25$ で、$30^2=900$ とほぼ同じです（0.77と0.64は四捨五入した値なので、少しずれる）。方眼紙に図を大きくかいてはかっても、同じ長さになります。

\bksubsection{3}{\ang{30}・\ang{45}・\ang{60}の三角比と、電卓}

\leavevmode\label{bk:4-6:atai}\ang{30}・\ang{45}・\ang{60}の三角比は、三角定規の辺の比（4-4）から、ぴったりの値が求められます。

\begin{bkfigure}
\begin{tikzpicture}[x=1mm, y=1mm]
  % 45°の三角定規 1:1:√2
  \coordinate (A) at (0,0); \coordinate (B) at (26,0); \coordinate (C) at (26,26);
  \fill[bkBlue!10] (A) -- (B) -- (C) -- cycle;
  \draw[bkBlue!80!black, line width=0.9pt] (A) -- (B) -- (C) -- cycle;
  \bkivgright{(B)}{90}{2.6}
  \draw[bkInk, line width=0.5pt] (7,0) arc (0:45:7);
  \node[font=\sffamily\footnotesize] at (22.5:10.5) {\ang{45}};
  \node[font=\sffamily\footnotesize, anchor=north] at (13,-1.5) {1};
  \node[font=\sffamily\footnotesize, anchor=west] at (27.5,13) {1};
  \node[font=\sffamily\footnotesize, anchor=south east] at (12,14) {$\sqrt2$};
  % 30°・60°の三角定規 1:2:√3
  \begin{scope}[shift={(58,0)}]
  \coordinate (P) at (0,0); \coordinate (Q) at ({20*sqrt(3)},0); \coordinate (R) at ({20*sqrt(3)},20);
  \fill[bkOrange!12] (P) -- (Q) -- (R) -- cycle;
  \draw[bkOrange!80!black, line width=0.9pt] (P) -- (Q) -- (R) -- cycle;
  \bkivgright{(Q)}{90}{2.6}
  \draw[bkInk, line width=0.5pt] (8,0) arc (0:30:8);
  \node[font=\sffamily\footnotesize] at (15:11.5) {\ang{30}};
  \draw[bkInk, line width=0.5pt] ($(R)+(210:6)$) arc (210:270:6);
  \node[font=\sffamily\footnotesize] at ($(R)+(240:9.5)$) {\ang{60}};
  \node[font=\sffamily\footnotesize, anchor=north] at ({10*sqrt(3)},-1.5) {$\sqrt3$};
  \node[font=\sffamily\footnotesize, anchor=west] at ({20*sqrt(3)+1.5},10) {1};
  \node[font=\sffamily\footnotesize, anchor=south east] at ({10*sqrt(3)-1},11) {2};
  \end{scope}
\end{tikzpicture}
\bkfigcaption{左は正方形を対角線で半分にした形。右は正三角形を半分にした形（4-4）。}
\end{bkfigure}

\noindent\textbf{\ang{45}}：左の三角形で、\ang{45}から見ると、対辺1、隣辺1、斜辺 $\sqrt2$ です。
\bkivgline{$\sin\ang{45}=\dfrac{1}{\sqrt2}=\dfrac{\sqrt2}{2}$\qquad $\cos\ang{45}=\dfrac{1}{\sqrt2}=\dfrac{\sqrt2}{2}$\qquad $\tan\ang{45}=\dfrac{1}{1}=1$}
$\dfrac{1}{\sqrt2}$ の分母と分子に $\sqrt2$ をかけて、分母を有理化しました（3-10）。

\noindent\textbf{\ang{30}}：右の三角形で、\ang{30}から見ると、対辺1、隣辺 $\sqrt3$、斜辺2です。
\bkivgline{$\sin\ang{30}=\dfrac{1}{2}$\qquad $\cos\ang{30}=\dfrac{\sqrt3}{2}$\qquad $\tan\ang{30}=\dfrac{1}{\sqrt3}=\dfrac{\sqrt3}{3}$}
\noindent\textbf{\ang{60}}：同じ三角形を、こんどは\ang{60}から見ます。\ang{60}の向かいの辺は $\sqrt3$ なので、対辺 $\sqrt3$、隣辺1、斜辺2です。
\bkivgline{$\sin\ang{60}=\dfrac{\sqrt3}{2}$\qquad $\cos\ang{60}=\dfrac{1}{2}$\qquad $\tan\ang{60}=\dfrac{\sqrt3}{1}=\sqrt3$}
同じ三角形でも、見る角が変わると、対辺と隣辺が入れかわります。

これから、≈ という記号を使います。≈ は「ほぼ等しい」という意味です（5-8 でくわしく）。$\sqrt2\approx1.414$、$\sqrt3\approx1.732$ なので（2-8）、小数にすると次の表になります。$\dfrac{\sqrt3}{2}\approx1.732\div2=0.866$ です。

\begin{bkcenter}
\begin{bktable}{colspec={Q[c,m,wd=18mm] Q[c,m,wd=30mm] Q[c,m,wd=30mm] Q[c,m,wd=30mm]}, rowsep=2pt, row{1}={font=\sffamily\small}}
$\theta$ & \ang{30} & \ang{45} & \ang{60} \\
$\sin\theta$ & $\dfrac{1}{2}=0.5$ & $\dfrac{\sqrt2}{2}\approx0.707$ & $\dfrac{\sqrt3}{2}\approx0.866$ \\
$\cos\theta$ & $\dfrac{\sqrt3}{2}\approx0.866$ & $\dfrac{\sqrt2}{2}\approx0.707$ & $\dfrac{1}{2}=0.5$ \\
$\tan\theta$ & $\dfrac{\sqrt3}{3}\approx0.577$ & 1 & $\sqrt3\approx1.732$ \\
\end{bktable}
\bktabcaption{\bkyattemiyouref ではかった\ang{30}の比（0.50、0.87、0.58）と合っている。}
\end{bkcenter}

\bkivghead{ほかの角は、電卓で}
\noindent\begin{minipage}[c]{0.56\linewidth}
\ang{40}や\ang{4.8}のような角の三角比は、関数電卓で求めます。手順は次のとおりです。
\begin{enumerate}[label={\bkcircled[bkOrange]{\arabic*}}, leftmargin=2.4\zw, itemsep=1pt, topsep=2pt]
\item 画面に「DEG」（度）と出ているか確かめる。「RAD」なら、度に切りかえる（4-2）。
\item \texttt{sin} のキーを押して、40と入れ、＝を押す。
\item 0.6427876… と出る（やってみようの0.64と合う）。
\end{enumerate}
\end{minipage}\hfill
\begin{minipage}[c]{0.4\linewidth}\centering
\begin{tikzpicture}[x=1mm, y=1mm]
  \bkdisplay{(0,14)}{46}{\footnotesize sin(40)\quad 0.6427876}
  \node[font=\sffamily\footnotesize\bfseries, text=bkInk, anchor=south west] at (1,25.6) {DEG};
  \bkdisplay{(0,0)}{46}{\footnotesize sin(30)\quad $-0.988$}
  \node[font=\sffamily\footnotesize\bfseries, text=bkRed!80!black, anchor=north west] at (1,-1.2) {RAD のまま（まちがい）};
\end{tikzpicture}
\end{minipage}

\par\vspace{4pt}
RAD のまま $\sin30$ を計算すると $-0.988$ になります。これは「30ラジアン」のサインで、\ang{30}のサインではありません。\ang{30}は $\dfrac{1}{2}=0.5$ と知っているので、$\sin30$ で0.5が出るかどうかで、モードを確かめられます。表計算なら \texttt{=SIN(RADIANS(40))} のように、度をラジアンに直してから計算します（4-2）。

\noindent 電卓で、\ang{10}ごとの値を表にしました。

\begin{bkcenter}
\begin{bktable}{colspec={Q[c,m,wd=12mm] *{9}{Q[c,m,wd=10.4mm]}}, rowsep=1.2pt, colsep=2pt, row{1}={font=\sffamily\footnotesize}, cells={font=\footnotesize}}
$\theta$ & \ang{10} & \ang{20} & \ang{30} & \ang{40} & \ang{50} & \ang{60} & \ang{70} & \ang{80} & \ang{90} \\
$\sin\theta$ & 0.174 & 0.342 & 0.500 & 0.643 & 0.766 & 0.866 & 0.940 & 0.985 & 1 \\
$\cos\theta$ & 0.985 & 0.940 & 0.866 & 0.766 & 0.643 & 0.500 & 0.342 & 0.174 & 0 \\
$\tan\theta$ & 0.176 & 0.364 & 0.577 & 0.839 & 1.192 & 1.732 & 2.747 & 5.671 & なし \\
\end{bktable}
\bktabcaption{小数第4位を四捨五入した値。$\sin$ の列を右から読むと、$\cos$ の列になっている（⑤）。}
\end{bkcenter}

\noindent\begin{minipage}[c]{0.5\linewidth}\centering
\begin{tikzpicture}
\begin{axis}[bkplot, width=56mm, height=44mm, xmin=0, xmax=98, ymin=0, ymax=1.18, xtick={30,60,90}, ytick={0.5,1}, xlabel={$\theta$}]
  \addplot[bkcurve, domain=0:90, samples=60] {sin(x)};
  \addplot[bkcurve2, domain=0:90, samples=60] {cos(x)};
  \node[font=\sffamily\footnotesize, text=bkBlue!80!black, anchor=west] at (axis cs:64,0.76) {$\sin\theta$};
  \node[font=\sffamily\footnotesize, text=bkOrange!80!black, anchor=west] at (axis cs:4,0.82) {$\cos\theta$};
\end{axis}
\end{tikzpicture}
\end{minipage}\hfill
\begin{minipage}[c]{0.46\linewidth}
$\theta$ が大きくなると、坂は急になります。だから高さの比 $\sin\theta$ は0から1へふえ、横の比 $\cos\theta$ は1から0へへります。斜辺がいちばん長い辺なので、$\sin\theta$ も $\cos\theta$ も1より大きくはなりません。$\tan\theta$ は、\ang{90}に近づくと、どこまでも大きくなります（横が0に近づくから）。\ang{90}の $\tan$ は決まりません。
\end{minipage}

\bksubsection{4}{$\tan\theta$ は坂の傾き}

\leavevmode\label{bk:4-6:tan}坂の傾きは、高さ÷横でした（3-5）。坂の角を $\theta$ とすると、高さは $\theta$ の対辺、横は隣辺です。だから、
\bkivgline{\fbox{傾き＝$\dfrac{\text{高さ}}{\text{横}}=\dfrac{\text{対辺}}{\text{隣辺}}=\tan\theta$}}
\textbf{$\tan\theta$ は、角 $\theta$ の坂の傾き}です。①の表の角度は、この関係で計算できます。電卓の表を使って確かめましょう。
\begin{itemize}[leftmargin=1.2\zw, itemsep=1pt, topsep=2pt, label={\color{\bkTT}●}]
\item 傾き1の坂：$\tan\ang{45}=1$ なので、角は\ang{45}。
\item 傾き2の坂：$\tan\ang{60}=1.732$、$\tan\ang{70}=2.747$ なので、角は\ang{60}と\ang{70}のあいだ。電卓で求めると\ang{63.4}です（⑦）。
\end{itemize}

\noindent\begin{minipage}[c]{0.56\linewidth}
グラフの直線でも同じです。原点を通る直線 $y=ax$（2-15）が、$x$ 軸の正の向きとつくる角を $\theta$ とします。直線の上の点 $(1,\ a)$ から $x$ 軸へ垂直な線を下ろすと、横1、たて $a$ の直角三角形ができます。だから
\bkivgline{$a=\dfrac{a}{1}=\tan\theta$}
\textbf{直線の傾き $a$ は、直線と $x$ 軸がつくる角の $\tan$} です。
\end{minipage}\hfill
\begin{minipage}[c]{0.4\linewidth}\centering
\begin{tikzpicture}
\begin{axis}[bkplot, width=50mm, height=44mm, xmin=-0.3, xmax=2.2, ymin=-0.2, ymax=1.75, xtick={1}, ytick=\empty, xlabel={$x$}, ylabel={$y$}]
  \addplot[bkcurve, domain=0:1.75] {0.8*x};
  \draw[bkOrange, line width=1.2pt] (axis cs:1,0) -- (axis cs:1,0.8);
  \addplot[only marks, mark=*, mark size=1.6pt, bkInk] coordinates {(1,0.8)};
  \node[font=\sffamily\footnotesize, anchor=north west] at (axis cs:1.06,0.76) {$(1,\ a)$};
  \node[font=\sffamily\footnotesize, text=bkOrange!80!black, anchor=east] at (axis cs:0.96,0.36) {$a$};
  \draw[bkInk, line width=0.5pt] (axis cs:0.3,0) arc[start angle=0, end angle=38.66, radius=8mm];
  \node[font=\sffamily\footnotesize] at (axis cs:0.48,0.13) {$\theta$};
  \node[font=\sffamily\footnotesize, text=bkBlue!80!black, anchor=north west] at (axis cs:1.62,1.18) {$y=ax$};
\end{axis}
\end{tikzpicture}
\end{minipage}

\bksubsection{5}{3つの三角比のつながり}

\leavevmode\label{bk:4-6:tsunagari}$\sin$・$\cos$・$\tan$ は、同じ三角形の辺の比なので、たがいにつながっています。

\bkivgstep{1}{$\tan\theta=\dfrac{\sin\theta}{\cos\theta}$}
$\sin\theta$ を $\cos\theta$ でわってみます。どちらも分母が斜辺なので、
\begin{align*}
\frac{\sin\theta}{\cos\theta}&=\frac{\text{対辺}}{\text{斜辺}}\div\frac{\text{隣辺}}{\text{斜辺}} && \text{\small（$\sin$ と $\cos$ の決め方を入れる）}\\
&=\frac{\text{対辺}}{\text{斜辺}}\times\frac{\text{斜辺}}{\text{隣辺}}=\frac{\text{対辺}}{\text{隣辺}}=\tan\theta && \text{\small（わり算は、逆数のかけ算。2-4）}
\end{align*}
斜辺どうしが約分されて消えます。数で確かめます。\ang{30}では、$0.5\div0.866$ を計算すると0.577になり、$\tan\ang{30}$ と合います。斜辺を1にした図で見れば、「たて $\sin\theta$ ÷ 横 $\cos\theta$」が傾き $\tan\theta$ です。

\bkivgstep{2}{$\sin^2\theta+\cos^2\theta=1$}
\noindent\begin{minipage}[c]{0.6\linewidth}
斜辺1の直角三角形（②の棒の図）で、三平方の定理（4-4）を使います。たてが $\sin\theta$、横が $\cos\theta$、斜辺が1なので、
\bkivgline{$(\sin\theta)^2+(\cos\theta)^2=1^2=1$}
$(\sin\theta)^2$ を、かっこを省いて $\sin^2\theta$ と書きます（$\sin$ の右上に2）。$\sin(\theta^2)$ のことではありません。
\end{minipage}\hfill
\begin{minipage}[c]{0.36\linewidth}\centering
\begin{tikzpicture}[x=1mm, y=1mm]
  \def\t{35}\def\L{36}
  \fill[bkTheme!12] (0,0) -- ({\L*cos(\t)},0) -- ({\L*cos(\t)},{\L*sin(\t)}) -- cycle;
  \draw[bkInk, line width=1.2pt] (0,0) -- ({\L*cos(\t)},{\L*sin(\t)});
  \draw[bkOrange, line width=1.4pt] ({\L*cos(\t)},0) -- ({\L*cos(\t)},{\L*sin(\t)});
  \draw[bkBlue, line width=1.4pt] (0,0) -- ({\L*cos(\t)},0);
  \bkivgright{({\L*cos(\t)},0)}{90}{2.2}
  \node[font=\sffamily\footnotesize, anchor=south east] at ({\L*cos(\t)/2-1},{\L*sin(\t)/2+1}) {1};
  \node[font=\sffamily\footnotesize, text=bkOrange!80!black, anchor=west] at ({\L*cos(\t)+1.5},{\L*sin(\t)/2}) {$\sin\theta$};
  \node[font=\sffamily\footnotesize, text=bkBlue!80!black, anchor=north] at ({\L*cos(\t)/2},-1.5) {$\cos\theta$};
  \draw[bkInk, line width=0.5pt] (6,0) arc (0:\t:6);
  \node[font=\sffamily\footnotesize] at ({\t/2}:9) {$\theta$};
\end{tikzpicture}
\end{minipage}
\bkivgline{\fbox{$\sin^2\theta+\cos^2\theta=1$}}
数で確かめます。
\bkivgline{\begin{tabular}{@{}l@{\quad}l@{}}\ang{30}： & $0.5^2+0.866^2=0.25+0.75=1$\\[2pt] \ang{40}： & $0.643^2+0.766^2=0.413+0.587=1.000$\end{tabular}}
合っています。この式を使うと、$\sin\theta$ がわかれば $\cos\theta$ がわかります。\bkreidairef{6}でためします。

\bkivgstep{3}{$\ang{90}-\theta$ から見ると、$\sin$ と $\cos$ が入れかわる}
\noindent\begin{minipage}[c]{0.6\linewidth}
直角三角形の直角でない2つの角は、$\theta$ と $\ang{90}-\theta$ です（内角の和が\ang{180}だから。4-1）。$\ang{90}-\theta$ の所に目を置くと、$\theta$ のときの対辺が隣辺に、隣辺が対辺に入れかわります。斜辺は同じです。だから、
\end{minipage}\hfill
\begin{minipage}[c]{0.36\linewidth}\centering
\begin{tikzpicture}[x=1mm, y=1mm]
  \coordinate (A) at (0,0); \coordinate (B) at (40,0); \coordinate (C) at (40,26);
  \fill[bkTheme!12] (A) -- (B) -- (C) -- cycle;
  \draw[bkInk!80, line width=0.9pt] (A) -- (B) -- (C) -- cycle;
  \draw[bkRed, line width=1.6pt] (B) -- (C);
  \draw[bkGreen!70!black, line width=1.6pt] (A) -- (B);
  \bkivgright{(B)}{90}{2.2}
  \draw[bkInk, line width=0.5pt] (7,0) arc (0:33:7);
  \node[font=\sffamily\footnotesize] at (16:10) {$\theta$};
  \draw[bkInk, line width=0.5pt] ($(C)+(213:6)$) arc (213:270:6);
  \node[font=\sffamily\footnotesize, anchor=east] at ($(C)+(-1.5,-1)$) {$\ang{90}-\theta$};
  \node[font=\sffamily\footnotesize, text=bkRed!80!black, anchor=west] at (41.5,13) {$a$};
  \node[font=\sffamily\footnotesize, text=bkGreen!50!black, anchor=north] at (20,-1.5) {$b$};
  \node[font=\sffamily\footnotesize, anchor=south east] at (19,14) {$c$};
\end{tikzpicture}
\end{minipage}
\bkivgline{$\sin(\ang{90}-\theta)=\dfrac{b}{c}=\cos\theta$\qquad $\cos(\ang{90}-\theta)=\dfrac{a}{c}=\sin\theta$\qquad $\tan(\ang{90}-\theta)=\dfrac{b}{a}=\dfrac{1}{\tan\theta}$}
③の電卓の表で、$\sin\ang{20}=0.342=\cos\ang{70}$ のように、$\sin$ の列を右から読むと $\cos$ の列になっていたのは、このためです（$\ang{20}+\ang{70}=\ang{90}$）。

\begin{motto}[2]
\textbf{「余弦」の「余」}\par
$\theta$ と、たして\ang{90}になる角 $\ang{90}-\theta$ を、$\theta$ の\textbf{余角}（よかく）といいます。上の式 $\cos\theta=\sin(\ang{90}-\theta)$ は、「$\theta$ のコサインは、余角のサイン」という意味です。余角（complement）のサイン（sine）なので、co-sine、日本語で余弦とよびます。$\tan(\ang{90}-\theta)=\dfrac{1}{\tan\theta}$ も同じ考え方で、$\dfrac{1}{\tan\theta}$ を cotangent（コタンジェント）ということがあります。
\end{motto}

\bksubsection{6}{斜面の重力を2つに分ける}

\leavevmode\label{bk:4-6:bunkai}いよいよ、\bkhatenaref の台車を考えます。

\bkivgstep{1}{重力の大きさ $mg$}
物を地球が引く力を\textbf{重力}といいます。質量約\qty{100}{g}の物にはたらく重力は、約\qty{1}{N}でした（2-15）。くわしくはかると、質量\qty{1}{kg}の物にはたらく重力は\qty{9.8}{N}です。9.8は、地球の上ではかって決まった値です（物が落ちるとき、速さが1秒ごとに約\qty{9.8}{m/s}ずつふえることと同じ数。第2巻の序章）。重力は質量に比例するので（2-15）、質量 $m$（kg）の物にはたらく重力は
\bkivgline{$9.8\times m=mg$（N）\qquad（$g=9.8$。1\,kgあたり9.8\,Nなので、$\qty{9.8}{N/kg}$ とも書く）}
です。$g$ は重力加速度ともよばれ、単位を $\unit{m/s^2}$ と書くこともあります（\unit{N/kg} と \unit{m/s^2} は同じ単位。5-5 でくわしく学びます）。質量\qty{20}{kg}の台車なら、重力は $20\times9.8=\qty{196}{N}$ です。質量（kg）は物の量で、月へ行っても変わりません。重力（N）は引っぱる力で、月では約 $\dfrac{1}{6}$ になります。

\bkivgstep{2}{1本の力を、2本の力に分けて考える}
\noindent\begin{minipage}[c]{0.58\linewidth}
重い荷物にロープをつけて、ななめ上へ引くと、荷物は前へ進みます。同時に、少し持ち上げられて軽くなります。ななめの1本の力が、「前へ引く力」と「上へ持ち上げる力」の2つの働きをしているのです。\par\vspace{2pt}
力は矢印で表します。力の矢印を\textbf{長方形の対角線}にすると、長方形の2つの辺が、分けた2つの力になります。この分け方で、もとの力と同じ働きになることは、実験で確かめられています（法則。4-1。4-10 でくわしく）。力を2つに分けることを、力の\textbf{分解}（ぶんかい）といいます。\bkidx{ちからのぶんかい}{力の分解}
\end{minipage}\hfill
\begin{minipage}[c]{0.38\linewidth}\centering
\begin{tikzpicture}[x=0.88mm, y=0.88mm]
  \fill[bkGray!30] (-4,-1.2) rectangle (50,0);
  \bkivgbox{(6,0)}{0}{bkOrange}
  \coordinate (O) at (12,5);
  \draw[bkGray, line width=0.5pt, dashed] (O) -- ++(30,0) -- ++(0,17.3);
  \draw[bkGray, line width=0.5pt, dashed] (O) -- ++(0,17.3) -- ++(30,0);
  \draw[bk arrow, bkInk] (O) -- ++(30,17.3);
  \draw[bk arrow, bkRed, line width=1.1pt] (O) -- ++(30,0);
  \draw[bk arrow, bkBlue, line width=1.1pt] (O) -- ++(0,17.3);
  \node[font=\sffamily\footnotesize, anchor=south east] at ($(O)+(13,9)$) {引く力};
  \node[font=\sffamily\footnotesize, text=bkRed!80!black, anchor=south] at ($(O)+(24,1)$) {前へ};
  \node[font=\sffamily\footnotesize, text=bkBlue!80!black, anchor=east] at ($(O)+(-1.2,15)$) {上へ};
\end{tikzpicture}
\end{minipage}

\bkivgstep{3}{斜面にそった向きと、斜面に垂直な向きに分ける}
斜面の上の台車では、重力 $mg$ を、次の2つの向きに分けます。
\begin{itemize}[leftmargin=1.2\zw, itemsep=1pt, topsep=2pt, label={\color{\bkTT}●}]
\item \textbf{斜面にそって下る向き}：台車を坂の下へ動かそうとする力。
\item \textbf{斜面に垂直で、斜面へ押しつける向き}：台車を斜面に押しつける力。
\end{itemize}
重力の矢印 $mg$ を対角線にして、この2つの向きに辺をもつ長方形をかきます。

\begin{bkfigure}
\begin{tikzpicture}[x=1mm, y=1mm]
  \def\t{30}
  \fill[bkGray!20] (0,0) -- (78,0) -- (78,{78*tan(\t)}) -- cycle;
  \draw[bkInk!80, line width=0.9pt] (0,0) -- (78,0) -- (78,{78*tan(\t)}) -- cycle;
  \draw[bkInk, line width=0.5pt] (12,0) arc (0:\t:12);
  \node[font=\sffamily\footnotesize] at ({\t/2}:15.5) {$\theta$};
  \coordinate (S) at ({57*cos(\t)},{57*sin(\t)});
  \bkivgcart{(S)}{\t}{bkOrange}
  \coordinate (G) at ($(S)+({\t+90}:6.8)$);
  \def\M{28}
  \coordinate (D) at ($(G)+(0,-\M)$);
  \coordinate (Pa) at ($(G)+({\t+180}:{\M*sin(\t)})$);
  \coordinate (Pp) at ($(G)+({\t-90}:{\M*cos(\t)})$);
  \fill[bkYellow!25] (G) -- (Pa) -- (D) -- (Pp) -- cycle;
  \draw[bkGray, line width=0.5pt, dashed] (Pa) -- (D) -- (Pp);
  \draw[bk arrow, bkInk, line width=1.5pt] (G) -- (D);
  \draw[bk arrow, bkRed, line width=1.4pt] (G) -- (Pa);
  \draw[bk arrow, bkBlue, line width=1.4pt] (G) -- (Pp);
  \fill[bkInk] (G) circle (0.8);
  \draw[bkInk, line width=0.6pt] ($(G)+(270:8)$) arc (270:{270+\t}:8);
  \node[font=\sffamily\footnotesize] at ($(G)+({270+\t/2}:11)$) {$\theta$};
  \bkivgright{(Pp)}{\t+90}{2.2}
  \node[font=\sffamily\small, anchor=west] at ($(D)+(1.5,0.5)$) {重力 $mg$};
  \node[font=\sffamily\small, text=bkRed!80!black, anchor=south east] at ($(Pa)+(-0.5,1)$) {$mg\sin\theta$};
  \node[font=\sffamily\small, text=bkBlue!80!black, anchor=west] at ($(Pp)+(1.6,0)$) {$mg\cos\theta$};
  \node[font=\sffamily\footnotesize, text=bkRed!80!black, anchor=west, align=left] at (84,42) {赤：斜面にそって\\下る向きの力};
  \node[font=\sffamily\footnotesize, text=bkBlue!80!black, anchor=west, align=left] at (84,29) {青：斜面へ\\押しつける力};
  \node[font=\sffamily\footnotesize, text=bkInk, anchor=west, align=left] at (84,16) {黒：重力\\（長方形の対角線）};
\end{tikzpicture}
\bkfigcaption{重力（黒）を対角線とする長方形の2辺が、分けた2つの力。重力と青の矢印のあいだの角が、斜面の角 $\theta$ と同じになる。}\label{bk:4-6:fig-shamen}
\end{bkfigure}

図の、黒（重力）・青・長方形の辺でできる直角三角形を見ます。斜辺は重力 $mg$ です。もし、黒と青のあいだの角が $\theta$ なら、$\theta$ から見て、青は隣辺、赤と同じ長さの辺（点線）は対辺です。だから、②の「かけ算の形」で
\bkivgline{\fbox{斜面にそった力＝$mg\sin\theta$}\qquad \fbox{斜面へ押しつける力＝$mg\cos\theta$}}
となります。のこる問題は、「黒と青のあいだの角は、ほんとうに斜面の角 $\theta$ と同じか」です。

\bkivgstep{4}{なぜ、斜面の角と同じ $\theta$ が出てくるのか}
\noindent\begin{minipage}[c]{0.56\linewidth}
\textbf{考え方1：三角形の角の和（4-1）}\par
荷物から真下へ線を引き、地面とぶつかる点をHとします。斜面・地面・この線で、大きな直角三角形ができます。角は、ふもとで $\theta$、Hで\ang{90}なので、荷物の所の角（斜面と真下の線のあいだ）は
\bkivgline{$\ang{180}-\ang{90}-\theta=\ang{90}-\theta$}
です。青の矢印は、斜面に垂直（\ang{90}）です。だから、真下の線（重力）と青のあいだの角は
\bkivgline{$\ang{90}-(\ang{90}-\theta)=\theta$}
となります。
\end{minipage}\hfill
\begin{minipage}[c]{0.4\linewidth}\centering
\begin{tikzpicture}[x=1mm, y=1mm]
  \def\t{30}
  \coordinate (S) at ({44*cos(\t)},{44*sin(\t)});
  \coordinate (H) at ({44*cos(\t)},0);
  \fill[bkGray!20] (0,0) -- (52,0) -- (52,{52*tan(\t)}) -- cycle;
  \draw[bkInk!80, line width=0.8pt] (0,0) -- (52,0) -- (52,{52*tan(\t)}) -- cycle;
  \fill[bkYellow!40] (0,0) -- (H) -- (S) -- cycle;
  \draw[bkInk, line width=0.9pt, dashed] (S) -- (H);
  \bkivgright{(H)}{90}{2.2}
  \node[font=\sffamily\footnotesize, anchor=north] at ($(H)+(0,-1.2)$) {H};
  \draw[bkInk, line width=0.5pt] (8,0) arc (0:\t:8);
  \node[font=\sffamily\footnotesize] at ({\t/2}:11) {$\theta$};
  \draw[bkInk, line width=0.5pt] ($(S)+(210:6)$) arc (210:270:6);
  \node[font=\sffamily\footnotesize] at ($(S)+(244:12.5)$) {$\ang{90}-\theta$};
  \draw[bk arrow, bkBlue, line width=1.1pt] (S) -- ++({\t-90}:16);
  \bkivgright{(S)}{\t-90}{2.2}
  \draw[bkInk, line width=0.6pt] ($(S)+(270:10)$) arc (270:{270+\t}:10);
  \node[font=\sffamily\footnotesize] at ($(S)+({270+\t/2}:13.2)$) {$\theta$};
  \fill[bkInk] (S) circle (0.8);
\end{tikzpicture}
\end{minipage}

\par\vspace{4pt}
\noindent\begin{minipage}[c]{0.56\linewidth}
\textbf{考え方2：床をかたむける}\par
はじめ、床は水平です。重力（真下）と、床に垂直な向きは、ぴったり重なっています。この床を $\theta$ だけかたむけて、斜面にします。床に垂直な向きも、床といっしょに $\theta$ だけ回ります。でも重力は、いつも真下のままです。だから、重力と「斜面に垂直な向き」のあいだには、床を回した角 $\theta$ だけずれができます。
\end{minipage}\hfill
\begin{minipage}[c]{0.4\linewidth}\centering
\begin{tikzpicture}[x=1mm, y=1mm]
  \def\t{30}
  \draw[bkGray, line width=1.2pt] (-22,0) -- (22,0);
  \draw[bkInk!80, line width=1.4pt] ({-22*cos(\t)},{-22*sin(\t)}) -- ({22*cos(\t)},{22*sin(\t)});
  \draw[bk arrow, bkInk, line width=1.1pt] (0,0) -- (0,-20);
  \draw[bk arrow, bkBlue, line width=1.1pt] (0,0) -- ({\t-90}:20);
  \draw[bkInk, line width=0.5pt] (16,0) arc (0:\t:16);
  \node[font=\sffamily\footnotesize] at ({\t/2}:19) {$\theta$};
  \draw[bkInk, line width=0.5pt] (0,-12) arc (270:{270+\t}:12);
  \node[font=\sffamily\footnotesize] at ({270+\t/2}:15) {$\theta$};
  \node[font=\sffamily\footnotesize, anchor=north] at (0,-21) {重力};
  \node[font=\sffamily\footnotesize, text=bkGray, anchor=south] at (-15,0.6) {水平な床};
\end{tikzpicture}
\end{minipage}

\bkivghead{考え方3：はしの場合で確かめる}
$\theta=\ang{0}$（平らな床）なら、$\sin\ang{0}=0$ で、斜面にそった力は0。台車は動きません。$\cos\ang{0}=1$ で、重力はすべて床へ押しつける力です。$\theta=\ang{90}$（かべ）なら、$\sin\ang{90}=1$ で、重力はすべて「かべにそって下る力」になり、物はまっすぐ落ちます。どちらも、ふだんの経験と合います。

\begin{seikaku}
\textbf{斜面が押し返す力}\par
台車が斜面にめりこまないのは、斜面が台車を、斜面に垂直な向きに押し返しているからです。この力を\textbf{垂直抗力}（すいちょくこうりょく）といいます。斜面に垂直な向きには台車は動かないので、垂直抗力の大きさは、押しつける力 $mg\cos\theta$ と同じです。物理シリーズでは、この力を $N$ と書きます。
\end{seikaku}

\bkivghead{\bkhatenaref の答え}
重力はまっすぐ下向きでも、斜面の上では、その一部 $mg\sin\theta$ が「斜面にそって下る力」として働くので、台車は坂を下ります。止めておくには、坂にそって上向きに、$mg\sin\theta$ の力で支えればよいのです（台車の車輪はよく回るので、まさつは考えません）。質量\qty{20}{kg}の台車なら、
\begin{itemize}[leftmargin=1.2\zw, itemsep=1pt, topsep=2pt, label={\color{\bkTT}●}]
\item \ang{30}の坂：$196\times\sin\ang{30}=196\times0.5=\qty{98}{N}$。重力の半分です。
\item 傾き $\dfrac{1}{12}$ のスロープ（約\ang{4.76}）：$196\times\sin\ang{4.76}\approx196\times0.083=\qty{16.3}{N}$。重力の約12分の1です（\bkreidairef{8}）。
\end{itemize}
ゆるやかな坂ほど $\sin\theta$ が小さいので、小さな力で止めておけます。スロープがゆるやかに作られているのは、このためです。

\bksubsection{7}{三角比から角度を求める：$\sin^{-1}$・$\cos^{-1}$・$\tan^{-1}$}

\leavevmode\label{bk:4-6:gyaku}④で、傾き2の坂の角は\ang{60}と\ang{70}のあいだ、とわかりました。ぴったりの角は、「$\tan\theta=2$ になる $\theta$」です。これを電卓で求めます。

\noindent\begin{minipage}[c]{0.56\linewidth}
関数電卓には、$\sin^{-1}$・$\cos^{-1}$・$\tan^{-1}$ のキーがあります（\texttt{SHIFT} を押してから $\sin$ などを押すことが多い）。$\tan^{-1}2$ は「$\tan$ が2になる角」という意味です。
\begin{enumerate}[label={\bkcircled[bkOrange]{\arabic*}}, leftmargin=2.4\zw, itemsep=1pt, topsep=2pt]
\item DEG になっているか確かめる。
\item $\tan^{-1}$ を押して、2と入れ、＝を押す。
\item 63.43… と出る。角は約\ang{63.4}。
\end{enumerate}
\end{minipage}\hfill
\begin{minipage}[c]{0.4\linewidth}\centering
\begin{tikzpicture}[x=1mm, y=1mm]
  \node[bk node, minimum width=20mm, minimum height=9mm, font=\small] (a) at (0,12) {角 \ang{63.4}};
  \node[bk node, minimum width=20mm, minimum height=9mm, font=\small] (b) at (0,-12) {傾き2};
  \draw[bk link] ([xshift=-4mm]a.south) -- ([xshift=-4mm]b.north);
  \draw[bk link] ([xshift=4mm]b.north) -- ([xshift=4mm]a.south);
  \node[bk link label, anchor=east] at (-4,0) {$\tan$};
  \node[bk link label, anchor=west] at (4,0) {$\tan^{-1}$};
\end{tikzpicture}
\end{minipage}

\par\vspace{4pt}
$\tan^{-1}$ は「アークタンジェント」と読みます。$\tan$ が「角→傾き」なら、$\tan^{-1}$ は反対の「傾き→角」です。右上の $-1$ は、$-1$ 乗（$\dfrac{1}{\tan}$）ではありません。「反対向きの計算」というしるしです。$\sin^{-1}$（アークサイン）、$\cos^{-1}$（アークコサイン）も同じです。表計算では \texttt{=DEGREES(ATAN(2))} のように書きます（\texttt{ATAN} の答えはラジアンなので、度に直す）。

これで、①の表の角度が計算できます。
\begin{bkcenter}
\begin{bktable}{colspec={Q[l,m,wd=34mm] Q[c,m,wd=20mm] Q[c,m,wd=46mm]}, rowsep=1.5pt, row{1}={font=\sffamily\small}}
 & 傾き & 角度 \\
スロープ & $\dfrac{1}{12}$ & $\tan^{-1}\dfrac{1}{12}\approx\ang{4.76}$ \\
道路の標識 & 0.1 & $\tan^{-1}0.1\approx\ang{5.71}$ \\
横1でたて2の坂 & 2 & $\tan^{-1}2\approx\ang{63.43}$ \\
\end{bktable}
\end{bkcenter}

\bksubsection{8}{\ang{0}から\ang{180}までの三角比}

\leavevmode\label{bk:4-6:donkaku}直角三角形の三角比は、\ang{90}より小さい角でしか決まりませんでした。でも、次の章（4-7）では、\ang{120}のように\ang{90}より大きい角（\textbf{鈍角}（どんかく））をもつ三角形を計算します。ロボットのうでのひじも、\ang{90}より大きく開きます。そこで、三角比の決め方を、\ang{180}まで広げます。

\bkivgstep{1}{半円の上の点で、決め直す}
\noindent\begin{minipage}[c]{0.52\linewidth}
原点Oを中心に、半径 $r$ の半円をかきます。$x$ 軸の正の向きから、左回りに角 $\theta$ だけ回った所にある半円の上の点を $\mathrm{P}(x,\ y)$ とします。このとき、
\bkivgline{$\sin\theta=\dfrac{y}{r}$\qquad $\cos\theta=\dfrac{x}{r}$\qquad $\tan\theta=\dfrac{y}{x}$}
と決めます。\par\vspace{2pt}
$\theta$ が鋭角のときは、Pから $x$ 軸へ垂直な線を下ろすと、斜辺 $r$、たて $y$、横 $x$ の直角三角形ができます。だから、この決め方は、②の決め方と同じ答えになります。
\end{minipage}\hfill
\begin{minipage}[c]{0.44\linewidth}\centering
\begin{tikzpicture}[x=0.85mm, y=0.85mm]
  \def\R{24}\def\t{40}
  \draw[bkGray, line width=0.6pt, -{Stealth[length=4pt]}] (-32,0) -- (32,0) node[anchor=north, font=\sffamily\footnotesize] {$x$};
  \draw[bkGray, line width=0.6pt, -{Stealth[length=4pt]}] (0,-3) -- (0,32) node[anchor=east, font=\sffamily\footnotesize] {$y$};
  \draw[bkBlue, line width=1.1pt] (\R,0) arc (0:180:\R);
  \fill[bkTheme!14] (0,0) -- ({\R*cos(\t)},0) -- ({\R*cos(\t)},{\R*sin(\t)}) -- cycle;
  \draw[bkInk, line width=0.9pt] (0,0) -- ({\R*cos(\t)},{\R*sin(\t)});
  \draw[bkOrange, line width=1.2pt] ({\R*cos(\t)},0) -- ({\R*cos(\t)},{\R*sin(\t)});
  \fill[bkInk] ({\R*cos(\t)},{\R*sin(\t)}) circle (0.9);
  \node[font=\sffamily\footnotesize, anchor=south west] at ({\R*cos(\t)+0.5},{\R*sin(\t)+0.3}) {$\mathrm{P}(x,\ y)$};
  \draw[bkInk, line width=0.5pt] (6,0) arc (0:\t:6);
  \node[font=\sffamily\footnotesize] at ({\t/2}:9) {$\theta$};
  \node[font=\sffamily\footnotesize, anchor=south east] at ({\R*cos(\t)/2},{\R*sin(\t)/2+0.5}) {$r$};
  \node[font=\sffamily\footnotesize, text=bkOrange!80!black, anchor=east] at ({\R*cos(\t)-0.8},{\R*sin(\t)/2+1}) {$y$};
  \node[font=\sffamily\footnotesize, anchor=north] at ({\R*cos(\t)/2},-1.2) {$x$};
  \node[font=\sffamily\footnotesize, anchor=north] at (0,-1.2) {O};
\end{tikzpicture}
\end{minipage}

\par\vspace{4pt}
この決め方は、②の「極座標を $(x,\ y)$ に直す式」$x=r\cos\theta$、$y=r\sin\theta$ を、\ang{180}まで使えるように広げたものでもあります（両辺を $r$ でわった形）。半径を1にすると、$r=1$ なので、Pの座標がそのまま $(\cos\theta,\ \sin\theta)$ になります。この見方は 4-8 で大活躍します。

\bkivgstep{2}{鈍角では、$\cos$ と $\tan$ が負になる}
\noindent\begin{minipage}[c]{0.44\linewidth}\centering
\begin{tikzpicture}[x=0.85mm, y=0.85mm]
  \def\R{24}\def\t{150}
  \draw[bkGray, line width=0.6pt, -{Stealth[length=4pt]}] (-32,0) -- (32,0) node[anchor=north, font=\sffamily\footnotesize] {$x$};
  \draw[bkGray, line width=0.6pt, -{Stealth[length=4pt]}] (0,-3) -- (0,32) node[anchor=east, font=\sffamily\footnotesize] {$y$};
  \draw[bkBlue, line width=1.1pt] (\R,0) arc (0:180:\R);
  % 150° の点 P と、30° の点 Q（y 軸について対称）
  \fill[bkOrange!15] (0,0) -- ({\R*cos(\t)},0) -- ({\R*cos(\t)},{\R*sin(\t)}) -- cycle;
  \fill[bkTheme!14] (0,0) -- ({\R*cos(30)},0) -- ({\R*cos(30)},{\R*sin(30)}) -- cycle;
  \draw[bkInk, line width=0.9pt] (0,0) -- ({\R*cos(\t)},{\R*sin(\t)});
  \draw[bkInk!50, line width=0.7pt] (0,0) -- ({\R*cos(30)},{\R*sin(30)});
  \draw[bkOrange, line width=1.2pt] ({\R*cos(\t)},0) -- ({\R*cos(\t)},{\R*sin(\t)});
  \draw[bkGray, line width=0.5pt, dashed] ({\R*cos(\t)},{\R*sin(\t)}) -- ({\R*cos(30)},{\R*sin(30)});
  \fill[bkInk] ({\R*cos(\t)},{\R*sin(\t)}) circle (0.9);
  \fill[bkInk!60] ({\R*cos(30)},{\R*sin(30)}) circle (0.8);
  \node[font=\sffamily\footnotesize, anchor=south east] at ({\R*cos(\t)-0.5},{\R*sin(\t)+0.3}) {P};
  \node[font=\sffamily\footnotesize, anchor=south west] at ({\R*cos(30)+0.5},{\R*sin(30)+0.3}) {Q};
  \draw[bkInk, line width=0.5pt] (5,0) arc (0:\t:5);
  \node[font=\sffamily\footnotesize] at (100:8.2) {$\theta$};
  \node[font=\sffamily\footnotesize, anchor=north] at ({\R*cos(\t)/2},-1.2) {$x<0$};
  \node[font=\sffamily\footnotesize, anchor=north] at (0,-1.2) {O};
\end{tikzpicture}
\end{minipage}\hfill
\begin{minipage}[c]{0.52\linewidth}
$\theta=\ang{150}$ の点Pは、$y$ 軸より左にあるので、$x$ 座標が負です。だから $\cos\ang{150}=\dfrac{x}{r}$ は負、$\tan\ang{150}=\dfrac{y}{x}$ も負です。$y$ 座標は正なので、$\sin\ang{150}$ は正です。\par\vspace{2pt}
Pを $y$ 軸について折り返すと、\ang{30}の点Qに重なります。PとQは、高さ $y$ が同じで、$x$ は符号だけが反対です。\ang{150}は $\ang{180}-\ang{30}$ なので、
\end{minipage}
\bkivgline{\fbox{\begin{tabular}{@{}l@{}}$\sin(\ang{180}-\theta)=\sin\theta$\qquad $\cos(\ang{180}-\theta)=-\cos\theta$\\[2pt] $\tan(\ang{180}-\theta)=-\tan\theta$\end{tabular}}}
たとえば、
\bkivgline{$\sin\ang{150}=\sin\ang{30}=\dfrac{1}{2}$\qquad $\cos\ang{150}=-\cos\ang{30}=-\dfrac{\sqrt3}{2}$}$\tan\theta=\dfrac{y}{x}$ は直線OPの傾きなので、鈍角では右下がりの直線になり、傾きが負になる（3-5）、と考えてもかまいません。

\bkivgstep{3}{\ang{0}・\ang{90}・\ang{180}}
Pが $x$ 軸の上（$\theta=\ang{0}$）なら $\mathrm{P}(r,\ 0)$、真上（\ang{90}）なら $(0,\ r)$、左はし（\ang{180}）なら $(-r,\ 0)$ です。表にまとめます。
\begin{bkcenter}
\begin{bktable}{colspec={Q[c,m,wd=12mm] *{9}{Q[c,m,wd=10.4mm]}}, rowsep=2pt, colsep=2pt, row{1}={font=\sffamily\footnotesize}}
$\theta$ & \ang{0} & \ang{30} & \ang{45} & \ang{60} & \ang{90} & \ang{120} & \ang{135} & \ang{150} & \ang{180} \\
$\sin\theta$ & 0 & $\frac{1}{2}$ & $\frac{\sqrt2}{2}$ & $\frac{\sqrt3}{2}$ & 1 & $\frac{\sqrt3}{2}$ & $\frac{\sqrt2}{2}$ & $\frac{1}{2}$ & 0 \\
$\cos\theta$ & 1 & $\frac{\sqrt3}{2}$ & $\frac{\sqrt2}{2}$ & $\frac{1}{2}$ & 0 & $-\frac{1}{2}$ & $-\frac{\sqrt2}{2}$ & $-\frac{\sqrt3}{2}$ & $-1$ \\
$\tan\theta$ & 0 & $\frac{\sqrt3}{3}$ & 1 & $\sqrt3$ & なし & $-\sqrt3$ & $-1$ & $-\frac{\sqrt3}{3}$ & 0 \\
\end{bktable}
\bktabcaption{\ang{90}では $x=0$ になり、0でわれないので $\tan\ang{90}$ は決まらない（1-14）。}
\end{bkcenter}
$\sin\theta$ は\ang{0}から\ang{180}まで0以上で、\ang{90}で最大の1です。そのため、$\sin\theta=\dfrac{1}{2}$ になる角は、\ang{30}と\ang{150}の2つあります（Pの高さが半径の半分になる点は、半円の上に2つあるから）。$\cos\theta$ は、1から $-1$ までへりつづけるので、$\cos\theta$ の値から角は1つに決まります。鈍角の三角比も、電卓でそのまま計算できます（$\cos150$ で $-0.866$）。

\begin{taisetsu}
三角比（$\theta$ から見た辺の比）：$\sin\theta=\dfrac{\text{対辺}}{\text{斜辺}}$\quad $\cos\theta=\dfrac{\text{隣辺}}{\text{斜辺}}$\quad $\tan\theta=\dfrac{\text{対辺}}{\text{隣辺}}$\par\vspace{3pt}
対辺＝斜辺×$\sin\theta$、隣辺＝斜辺×$\cos\theta$。坂の傾き＝$\tan\theta$。\par\vspace{3pt}
$\tan\theta=\dfrac{\sin\theta}{\cos\theta}$\quad $\sin^2\theta+\cos^2\theta=1$\quad $\sin(\ang{90}-\theta)=\cos\theta$\par\vspace{3pt}
斜面の重力：斜面にそって $mg\sin\theta$、斜面へ押しつける向きに $mg\cos\theta$。\par\vspace{3pt}
極座標 $(r,\ \theta)$ を $(x,\ y)$ に直す：$x=r\cos\theta$、$y=r\sin\theta$\par\vspace{3pt}
$\sin(\ang{180}-\theta)=\sin\theta$\quad $\cos(\ang{180}-\theta)=-\cos\theta$（鈍角の $\cos$ は負）
\end{taisetsu}

\begin{kotoba}
\small
\bkword{斜辺・対辺・隣辺}{直角三角形の、直角の向かいの辺・角 $\theta$ の向かいの辺・角 $\theta$ のとなりの辺（斜辺でないほう）。}
\bkword{三角比}{$\sin\theta$（正弦）・$\cos\theta$（余弦）・$\tan\theta$（正接）のこと。角 $\theta$ だけで決まる。}
\bkword{こう配}{傾きを百分率（\%）で表したもの。傾き0.1はこう配10\%。}
\bkword{力の分解}{1つの力を、同じ働きをする2つの力に分けること。}
\bkword{垂直抗力}{面が物を、面に垂直な向きに押し返す力。}
\bkword{鋭角・鈍角}{\ang{0}より大きく\ang{90}より小さい角・\ang{90}より大きく\ang{180}より小さい角。}
\bkword{$\sin^{-1}$・$\tan^{-1}$}{三角比の値から、角を求める計算（アークサイン・アークタンジェント）。}
\end{kotoba}

\bkkigocards{\bkkigocard{\shortstack{$\sin$\\[2pt]$\cos$}}{サイン・コサイン}{角 $\theta$ の直角三角形の、斜辺に対する対辺・隣辺の比（正弦・余弦）}{sine, cosine}{4-6}}
{\bkkigocard{$\tan$}{タンジェント}{対辺÷隣辺の比（正接）。角 $\theta$ の坂の傾き}{tangent}{4-6}}

\begin{column}{ピラミッドの「セケド」}
今から\num{4500}年ほど前、エジプトの人々は、ピラミッドの面の急さを「セケド」という数で決めていました。セケドは「高さを1キュービット（約\qty{52}{cm}）上がるあいだに、横へ何パーム進むか」で、1キュービットは7パームです。ギザの大ピラミッドのセケドは $5\dfrac{1}{2}$ パームといわれます。高さ7に対して横5.5なので、傾きは $\dfrac{7}{5.5}\approx1.27$、角度は $\tan^{-1}\dfrac{7}{5.5}\approx\ang{51.8}$ です。実際にはかった大ピラミッドの面の角度も約\ang{51.8}です。4500年前の人も、「たて÷横」で坂を作っていたのです。
\end{column}

\begin{column}{スマートフォンは、重力の分け方でかたむきを知る}
スマートフォンには、加速度センサーという、とても小さな半導体の部品が入っています。画面を水平から $\theta$ だけかたむけると、センサーは、重力 $mg$ が「画面にそった向き」に $mg\sin\theta$、「画面に垂直な向き」に $mg\cos\theta$ だけ分かれているのをはかります（⑥と同じ分け方）。2つの値の比は
\bkivgline{$\dfrac{mg\sin\theta}{mg\cos\theta}=\dfrac{\sin\theta}{\cos\theta}=\tan\theta$}
なので、$\tan^{-1}$ で角度がわかります。質量\qty{1}{kg}あたりに直して、そった向きが\qty{4.9}{N}、垂直な向きが\qty{8.49}{N}なら、$\tan\theta=\dfrac{4.9}{8.49}\approx0.577$ で、$\theta\approx\ang{30}$ です。画面を回すと表示が横向きになるのも、水平器のアプリが坂の角度を表示するのも、この計算のおかげです。
\end{column}


%% ---------------- 例題 ----------------
\bksection{reidai}{例題でたしかめよう}

\begin{reidai}[直角三角形の三角比]
\noindent\begin{minipage}[t]{0.66\linewidth}
右の直角三角形ABC（Cが直角）で、$\sin A$、$\cos A$、$\tan A$ と、$\sin B$、$\cos B$、$\tan B$ を求めましょう。$\sin A$ は「角Aの $\sin$」のことです。\par
\kangaekata どの角から見るかで、対辺と隣辺が変わります。角Aに目を置くと、向かいの辺はBC、となりの辺はAC。角Bに目を置くと、向かいの辺はAC、となりの辺はBCです。斜辺は、どちらから見ても直角の向かいのABです。
\end{minipage}\hfill
\begin{minipage}[t]{0.3\linewidth}\vspace{0pt}\centering
\begin{tikzpicture}[x=0.8mm, y=0.8mm]
  \coordinate (A) at (0,0); \coordinate (C) at (32,0); \coordinate (B) at (32,24);
  \fill[bkBlue!10] (A) -- (C) -- (B) -- cycle;
  \draw[bkInk, line width=0.9pt] (A) -- (C) -- (B) -- cycle;
  \bkivgright{(C)}{90}{2.4}
  \node[font=\sffamily\footnotesize, anchor=east] at (-0.8,0) {A};
  \node[font=\sffamily\footnotesize, anchor=north west] at (32.5,-0.3) {C};
  \node[font=\sffamily\footnotesize, anchor=south] at (32,24.8) {B};
  \node[font=\sffamily\footnotesize, anchor=north] at (16,-1.4) {4};
  \node[font=\sffamily\footnotesize, anchor=west] at (33.4,12) {3};
  \node[font=\sffamily\footnotesize, anchor=south east] at (15.5,12.5) {5};
\end{tikzpicture}
\end{minipage}\par
\shiki 角Aから見る：対辺 $\mathrm{BC}=3$、隣辺 $\mathrm{AC}=4$、斜辺 $\mathrm{AB}=5$\par
\hspace*{2.6\zw}$\sin A=\dfrac{3}{5}$、$\cos A=\dfrac{4}{5}$、$\tan A=\dfrac{3}{4}$\par
\hspace*{2.6\zw}角Bから見る：対辺 $\mathrm{AC}=4$、隣辺 $\mathrm{BC}=3$、斜辺 $\mathrm{AB}=5$\par
\hspace*{2.6\zw}$\sin B=\dfrac{4}{5}$、$\cos B=\dfrac{3}{5}$、$\tan B=\dfrac{4}{3}$\par
\kotae $\sin A=\dfrac{3}{5}$、$\cos A=\dfrac{4}{5}$、$\tan A=\dfrac{3}{4}$、$\sin B=\dfrac{4}{5}$、$\cos B=\dfrac{3}{5}$、$\tan B=\dfrac{4}{3}$
\end{reidai}
$\sin A=\cos B$、$\cos A=\sin B$ になっています。$A+B=\ang{90}$ なので、⑤の「$\ang{90}-\theta$ から見ると入れかわる」のとおりです。

\begin{jibun}
\noindent\begin{minipage}[t]{0.68\linewidth}
右の直角三角形DEF（Fが直角）で、$\sin D$、$\cos D$、$\tan D$ と、$\sin E$、$\cos E$、$\tan E$ を求めましょう。
\end{minipage}\hfill
\begin{minipage}[t]{0.28\linewidth}\vspace{0pt}\centering
\begin{tikzpicture}[x=0.72mm, y=0.72mm]
  \coordinate (D) at (0,0); \coordinate (F) at (31.2,0); \coordinate (E) at (31.2,13);
  \draw[bkInk, line width=0.9pt] (D) -- (F) -- (E) -- cycle;
  \bkivgright{(F)}{90}{2.2}
  \node[font=\sffamily\footnotesize, anchor=east] at (-0.8,0) {D};
  \node[font=\sffamily\footnotesize, anchor=north west] at (31.7,-0.3) {F};
  \node[font=\sffamily\footnotesize, anchor=south] at (31.2,13.8) {E};
  \node[font=\sffamily\footnotesize, anchor=north] at (15.6,-1.4) {12};
  \node[font=\sffamily\footnotesize, anchor=west] at (32.6,6.5) {5};
  \node[font=\sffamily\footnotesize, anchor=south east] at (15.6,7) {13};
\end{tikzpicture}
\end{minipage}
\bkjans{$\sin D=\dfrac{5}{13}$、$\cos D=\dfrac{12}{13}$、$\tan D=\dfrac{5}{12}$、$\sin E=\dfrac{12}{13}$、$\cos E=\dfrac{5}{13}$、$\tan E=\dfrac{12}{5}$}
\end{jibun}
\textbf{ポイント}\quad 辺の長さが2つしかわからないときは、三平方の定理（4-4）でのこりの辺を求めてから、三角比を作ります。どの角から見るかで、対辺と隣辺が入れかわることも、もう一度確かめましょう。

\begin{reidai}[斜辺と角から、辺の長さを求める]
（1）長さ\qty{4}{m}のはしごを、地面と\ang{60}の角でかべに立てかけます。はしごの上のはしの高さと、はしごの下のはしからかべまでの距離を求めましょう。\par
（2）高さ\qty{1.5}{m}の台に、\ang{30}の角で板をかけてすべり台を作ります。板の長さは何mですか。\par
\kangaekata はしごが斜辺、高さが対辺、かべまでの距離が隣辺です。（1）は「対辺＝斜辺×$\sin\theta$」「隣辺＝斜辺×$\cos\theta$」。（2）は対辺がわかっていて斜辺を求めるので、$\sin\theta=\dfrac{\text{対辺}}{\text{斜辺}}$ を斜辺について解きます。\par
\shiki（1）高さ $=4\times\sin\ang{60}=4\times\dfrac{\sqrt3}{2}=2\sqrt3\approx2\times1.732=3.464$\par
\hspace*{4.6\zw}距離 $=4\times\cos\ang{60}=4\times\dfrac{1}{2}=2$\par
\hspace*{2.6\zw}（2）板の長さを $L$ とすると、$\sin\ang{30}=\dfrac{1.5}{L}$\par
\hspace*{4.6\zw}両辺に $L$ をかけて $L\sin\ang{30}=1.5$。両辺を $\sin\ang{30}$ でわって\par
\hspace*{4.6\zw}$L=\dfrac{1.5}{\sin\ang{30}}=\dfrac{1.5}{0.5}=3$\par
\kotae （1）高さ約\qty{3.46}{m}、距離\qty{2}{m}\quad（2）\qty{3}{m}
\end{reidai}
（2）は、「\ang{30}の坂では、高さは斜辺の半分」と考えても、斜辺は高さの2倍で\qty{3}{m}とわかります。

\begin{jibun}
（1）長さ\qty{200}{m}の坂道が、水平と\ang{5}の角をなしています。坂をのぼりきると何m高くなりますか。また、横に何m進みますか。電卓を使い、小数第1位まで求めましょう。\par
（2）高さ\qty{2}{m}の所まで、\ang{20}の角でまっすぐなスロープをかけます。スロープの長さは何mですか（小数第1位まで）。
\bkjans{（1）\nobreak 高さ $200\sin\ang{5}\approx\qty{17.4}{m}$、横 $200\cos\ang{5}\approx\qty{199.2}{m}$。（2）\nobreak $\dfrac{2}{\sin\ang{20}}\approx\qty{5.8}{m}$。}
\end{jibun}
\textbf{ポイント}\quad 斜辺がわかっていて、たてや横を求めるときは「かけ算」（斜辺×$\sin\theta$）、たてがわかっていて斜辺を求めるときは「わり算」（たて÷$\sin\theta$）です。どちらも、もとの式 $\sin\theta=\dfrac{\text{対辺}}{\text{斜辺}}$ から作れます。

\begin{reidai}[極座標を $(x,\ y)$ に直す]
（1）灯台のレーダーに、船Aが極座標 $(3,\ \ang{60})$ の所に見えました（距離の単位は km。4-2）。灯台を原点、東を $x$ 軸の正の向き、北を $y$ 軸の正の向きとして、船Aの $(x,\ y)$ を求めましょう。\par
（2）長さ\qty{50}{cm}のロボットのうでを、肩から\ang{25}の向きにのばしました。肩を原点として、手の $(x,\ y)$ を電卓で求めましょう（小数第1位まで）。\par
\kangaekata 斜辺 $r$ の直角三角形で、横が $x=r\cos\theta$、たてが $y=r\sin\theta$ です（②）。（1）は三角定規の値（③）、（2）は電卓（DEG）で計算します。\par
\shiki（1）$x=3\cos\ang{60}=3\times\dfrac{1}{2}=\dfrac{3}{2}=1.5$\par
\hspace*{4.6\zw}$y=3\sin\ang{60}=3\times\dfrac{\sqrt3}{2}=\dfrac{3\sqrt3}{2}\approx\dfrac{3\times1.732}{2}=2.598$\par
\hspace*{2.6\zw}（2）$x=50\cos\ang{25}=50\times0.9063=45.31\cdots$\par
\hspace*{4.6\zw}$y=50\sin\ang{25}=50\times0.4226=21.13\cdots$\par
\kotae （1）$\left(\dfrac{3}{2},\ \dfrac{3\sqrt3}{2}\right)$（約 $(1.5,\ 2.6)$）\quad（2）約 $(45.3,\ 21.1)$
\end{reidai}
（1）の答えは、4-2で方眼に図をかいて読みとった「右へ約1.5、上へ約2.6」と同じです。図をかかなくても、計算で求められました。

\begin{jibun}
（1）極座標が $(4,\ \ang{30})$ の点の $(x,\ y)$ を求めましょう（$\sqrt{\phantom{3}}$ を使って）。\par
（2）長さ\qty{40}{cm}のロボットのうでを、肩から\ang{55}の向きにのばしました。手の $(x,\ y)$ を、電卓で小数第1位まで求めましょう。
\bkjans{（1）\nobreak $x=4\cos\ang{30}=2\sqrt3$、$y=4\sin\ang{30}=2$ で、$(2\sqrt3,\ 2)$（約 $(3.46,\ 2)$）。（2）\nobreak $x=40\cos\ang{55}\approx22.9$、$y=40\sin\ang{55}\approx32.8$ で、約 $(22.9,\ 32.8)$。}
\end{jibun}
\textbf{ポイント}\quad $(r,\ \theta)$ から $(x,\ y)$ へは、「横は $r\cos\theta$、たては $r\sin\theta$」。答えは、$x^2+y^2=r^2$（三平方の定理）で確かめられます。（1）なら $1.5^2+2.598^2\approx9=3^2$ です。

\begin{reidai}[$\tan$ で高さをはかる]
\noindent\begin{minipage}[t]{0.6\linewidth}
校庭の木から\qty{20}{m}はなれた所に立ち、木のてっぺんを見上げると、水平から\ang{35}の角でした。目の高さは\qty{1.5}{m}です。木の高さは何mですか（小数第1位まで）。\par
\kangaekata 目の高さの水平な線と、木でできる直角三角形を考えます。\ang{35}から見て、隣辺が\qty{20}{m}、求めたい木の部分（目より上）が対辺です。「対辺＝隣辺×$\tan\theta$」を使い、最後に目の高さをたします。
\end{minipage}\hfill
\begin{minipage}[t]{0.36\linewidth}\vspace{0pt}\centering
\begin{tikzpicture}[x=0.8mm, y=0.8mm]
  \fill[bkGray!30] (-4,-1) rectangle (44,0);
  \bkkid{(0,0)}{8}{bkBlue}{bkInk}
  \coordinate (E) at (0,6.6);
  \coordinate (T) at (36,{6.6+36*tan(35)});
  \fill[bkStem] (35,0) rectangle (37.4,22);
  \fill[bkLeaf] (36.2,25.8) circle (6);
  \fill[bkLeaf!85!white] (32,21.5) circle (4) (40.4,21.5) circle (4);
  \fill[bkInk] (T) circle (0.7);
  \draw[bkInk, line width=0.6pt, dashed] (E) -- (36,6.6);
  \draw[bkOrange, line width=1pt] (E) -- (T);
  \draw[bkInk, line width=0.5pt] ($(E)+(9,0)$) arc (0:35:9);
  \node[font=\sffamily\footnotesize] at ($(E)+(17.5:12.5)$) {\ang{35}};
  \node[font=\sffamily\footnotesize, anchor=north] at (18,-1.6) {\qty{20}{m}};
  \draw[bkGray, line width=0.4pt, <->] (-4,0) -- (-4,6.6);
  \node[font=\sffamily\footnotesize, anchor=east] at (-4.6,3.3) {1.5};
\end{tikzpicture}
\end{minipage}\par
\shiki 目より上の部分 $=20\times\tan\ang{35}$\par
\hspace*{4.6\zw}$=20\times0.7002=14.004$\par
\hspace*{2.6\zw}木の高さ $=14.004+1.5=15.504$\par
\kotae 約\qty{15.5}{m}
\end{reidai}
見上げる角を\textbf{仰角}（ぎょうかく）といいます。仰角と距離がわかれば、のぼらなくても高さがわかります。目の高さをたしわすれないようにしましょう。

\begin{jibun}
ビルから\qty{30}{m}はなれた所で、屋上を見上げると、仰角は\ang{40}でした。目の高さを\qty{1.6}{m}として、ビルの高さを小数第1位まで求めましょう。
\bkjans{$30\tan\ang{40}+1.6\approx25.17+1.6$ で約\qty{26.8}{m}。}
\end{jibun}
\textbf{ポイント}\quad $\tan$ は「横がわかっているとき、たてを出す」道具です。仰角が\ang{45}なら、$\tan\ang{45}=1$ なので、目より上の高さは、はなれた距離と同じになります。

\begin{column}{ストローと分度器で、仰角をはかろう}
\noindent\begin{minipage}[c]{0.6\linewidth}
分度器のまっすぐな辺にストローをはり、中心の穴から糸をたらして、先にクリップなどのおもりをつけます。ストローで木のてっぺんをのぞくと、糸は分度器の\ang{90}の目もりから少しずれた所を指します。このずれが仰角です。\par\vspace{2pt}
わけは、⑥の「床をかたむける」と同じです。ストローを水平から $\theta$ だけ上に向けると、分度器の\ang{90}の線も $\theta$ だけ回ります。でも糸は、重力でいつも真下を向いています。だから、糸と\ang{90}の線のずれが $\theta$ になるのです。
\end{minipage}\hfill
\begin{minipage}[c]{0.36\linewidth}\centering
\begin{tikzpicture}[x=1mm, y=1mm]
  \def\t{30}
  \begin{scope}[rotate=\t]
    \fill[bkBlue!8] (-16,0) arc (180:360:16) -- cycle;
    \draw[bkBlue!70!black, line width=0.8pt] (-16,0) arc (180:360:16) -- cycle;
    \foreach \d in {190,200,...,350} {\draw[bkInk, line width=0.4pt] (\d:16) -- (\d:14);}
    \draw[bkRed, line width=0.8pt] (0,0) -- (270:16);
    \fill[bkOrange] (-20,0.4) rectangle (20,2.2);
    \draw[bkOrange!70!black, line width=0.5pt] (-20,0.4) rectangle (20,2.2);
  \end{scope}
  \draw[bkInk, line width=0.7pt] (0,0) -- (0,-22);
  \fill[bkGray] (0,-23.2) circle (1.4);
  \fill[bkInk] (0,0) circle (0.7);
  \draw[bkInk, line width=0.5pt] (0,-10) arc (270:{270+\t}:10);
  \node[font=\sffamily\footnotesize] at ({270+\t/2}:13) {$\theta$};
  \node[font=\sffamily\footnotesize, text=bkRed!80!black, anchor=west] at ($({270+\t}:16)+(0.8,-0.5)$) {\ang{90}の線};
  \node[font=\sffamily\footnotesize, anchor=east] at (-1.5,-20) {糸};
  \node[font=\sffamily\footnotesize, text=bkOrange!80!black, anchor=south] at ($(\t:20)+(-1,2.5)$) {ストロー};
\end{tikzpicture}
\end{minipage}
\end{column}

\begin{bkfigure}
\begin{tikzpicture}[x=1mm, y=1mm]
  \foreach \xx/\cA/\cB/\cC/\key/\lab in {0/bkRed/bkGray/bkBlue/{$\sin^{-1}$}/{対辺と斜辺}, 44/bkGray/bkGreen!70!black/bkBlue/{$\cos^{-1}$}/{隣辺と斜辺}, 88/bkRed/bkGreen!70!black/bkGray/{$\tan^{-1}$}/{対辺と隣辺}} {
    \begin{scope}[shift={(\xx,0)}]
      \fill[bkTheme!10] (0,0) -- (34,0) -- (34,18) -- cycle;
      \draw[\cA, line width=1.5pt] (34,0) -- (34,18);
      \draw[\cB, line width=1.5pt] (0,0) -- (34,0);
      \draw[\cC, line width=1.5pt] (0,0) -- (34,18);
      \draw[bkInk, line width=0.4pt] (31.8,0) -- (31.8,2.2) -- (34,2.2);
      \draw[bkInk, line width=0.5pt] (7,0) arc (0:27.9:7);
      \node[font=\sffamily\footnotesize] at (14:10) {$\theta$};
      \node[font=\sffamily\footnotesize, anchor=north] at (17,-1.5) {\lab};
      \node[font=\sffamily\small, anchor=north] at (17,-7) {\key};
    \end{scope}}
\end{tikzpicture}
\bkfigcaption{色のついた2辺がわかっているとき、使うキー。赤は対辺、緑は隣辺、青は斜辺。灰色の辺は使わない。}
\end{bkfigure}

\begin{reidai}[三角比から角度を求める]
（1）傾き $\dfrac{1}{12}$ のスロープの角度を求めましょう。\par
（2）長さ\qty{3}{m}のはしごを立てかけると、上のはしが高さ\qty{2.8}{m}になりました。はしごと地面の角度を求めましょう。\par
（3）長さ\qty{4}{m}のはしごの下のはしを、かべから\qty{1}{m}はなして立てかけました。はしごと地面の角度を求めましょう。\par
どれも、小数第1位まで求めます。\par
\kangaekata （1）傾き＝$\tan\theta$ なので、$\tan^{-1}$ で角を出します。（2）斜辺3と対辺2.8がわかっているので、$\sin\theta=\dfrac{2.8}{3}$。$\sin^{-1}$ で角を出します。（3）斜辺4と隣辺1（かべまでの距離）がわかっているので、$\cos\theta=\dfrac{1}{4}$。$\cos^{-1}$ で角を出します。電卓は DEG にします。\par
\shiki（1）$\tan\theta=\dfrac{1}{12}=0.0833\cdots$\par
\hspace*{4.6\zw}$\theta=\tan^{-1}\dfrac{1}{12}=4.763\cdots$\par
\hspace*{2.6\zw}（2）$\sin\theta=\dfrac{2.8}{3}=0.9333\cdots$\par
\hspace*{4.6\zw}$\theta=\sin^{-1}0.9333\cdots=68.96\cdots$\par
\hspace*{2.6\zw}（3）$\cos\theta=\dfrac{1}{4}=0.25$\par
\hspace*{4.6\zw}$\theta=\cos^{-1}0.25=75.52\cdots$\par
\kotae （1）約\ang{4.8}\quad（2）約\ang{69.0}\quad（3）約\ang{75.5}
\end{reidai}
（2）は、電卓に $2.8\div3$ を計算させてから $\sin^{-1}$ を押すと、途中で数を丸めずにすみます。

\begin{jibun}
（1）こう配10\%の道路の角度を求めましょう。\par
（2）長さ\qty{5}{m}のはしごの下のはしを、かべから\qty{1.5}{m}はなして立てかけました。はしごと地面の角度を求めましょう。\par
どちらも小数第1位まで求めましょう。
\bkjans{（1）\nobreak $\tan^{-1}0.1\approx\ang{5.7}$。（2）\nobreak $\cos\theta=\dfrac{1.5}{5}=0.3$ より約\ang{72.5}。}
\end{jibun}
\textbf{ポイント}\quad わかっている2辺が「対辺と斜辺」なら $\sin^{-1}$、「隣辺と斜辺」なら $\cos^{-1}$、「対辺と隣辺」なら $\tan^{-1}$ を使います。出てきた角を電卓で $\sin$ などに入れ直して、もとの比にもどるか確かめましょう。

\begin{reidai}[1つの三角比から、ほかの三角比を求める]
$\theta$ が鋭角で、$\sin\theta=\dfrac{3}{5}$ のとき、$\cos\theta$ と $\tan\theta$ を求めましょう。\par
\kangaekata $\sin^2\theta+\cos^2\theta=1$ に $\sin\theta$ の値を入れて、$\cos\theta$ を求めます。鋭角なので $\cos\theta$ は正です。$\tan\theta$ は $\dfrac{\sin\theta}{\cos\theta}$ です。\par
\shiki $\cos^2\theta=1-\sin^2\theta=1-\Bigl(\dfrac{3}{5}\Bigr)^2=1-\dfrac{9}{25}=\dfrac{16}{25}$\par
\hspace*{2.6\zw}$\cos\theta>0$ なので\quad $\cos\theta=\sqrt{\dfrac{16}{25}}=\dfrac{4}{5}$\par
\hspace*{2.6\zw}$\tan\theta=\dfrac{\sin\theta}{\cos\theta}=\dfrac{3}{5}\div\dfrac{4}{5}=\dfrac{3}{5}\times\dfrac{5}{4}=\dfrac{3}{4}$\par
\kotae $\cos\theta=\dfrac{4}{5}$、$\tan\theta=\dfrac{3}{4}$
\end{reidai}
図でも確かめられます。斜辺5、対辺3の直角三角形をかくと、隣辺は三平方の定理で $\sqrt{25-9}=4$。$\cos\theta=\dfrac{4}{5}$、$\tan\theta=\dfrac{3}{4}$ です（\bkreidairef{1}の三角形と同じ）。

\begin{jibun}
$\theta$ が鋭角で、$\cos\theta=\dfrac{5}{13}$ のとき、$\sin\theta$ と $\tan\theta$ を求めましょう。
\bkjans{$\sin^2\theta=\dfrac{144}{169}$ より $\sin\theta=\dfrac{12}{13}$、$\tan\theta=\dfrac{12}{5}$。}
\end{jibun}
\textbf{ポイント}\quad $\theta$ が鈍角なら、$\cos\theta$ は負になります。$\sqrt{\phantom{2}}$ をとったあと、角が鋭角か鈍角かで、符号を決めます。問題に「鋭角」と書いてあるのは、そのためです。

\begin{reidai}[\ang{0}〜\ang{180}の三角比]
（1）次の値を求めましょう。\par
\hspace*{2\zw}ア\enspace $\sin\ang{60}$\quad イ\enspace $\tan\ang{45}$\quad ウ\enspace $\sin\ang{135}$\quad エ\enspace $\cos\ang{135}$\quad オ\enspace $\tan\ang{150}$\par
（2）$\ang{0}\leqq\theta\leqq\ang{180}$ で、$\sin\theta=\dfrac{\sqrt3}{2}$ になる $\theta$ を求めましょう（$\leqq$ は「以下」）。\par
\kangaekata 鈍角は $\ang{180}-\theta$ の式で鋭角に直します。$\sin$ はそのまま、$\cos$ と $\tan$ は符号が反対になります。（2）は、半円の上で高さが同じ点が2つあることに注意します。\par
\shiki（1）ア\enspace $\dfrac{\sqrt3}{2}$\quad イ\enspace 1\par
\hspace*{4.6\zw}ウ\enspace $\sin\ang{135}=\sin(\ang{180}-\ang{45})=\sin\ang{45}=\dfrac{\sqrt2}{2}$\par
\hspace*{4.6\zw}エ\enspace $\cos\ang{135}=-\cos\ang{45}=-\dfrac{\sqrt2}{2}$\par
\hspace*{4.6\zw}オ\enspace $\tan\ang{150}=-\tan\ang{30}=-\dfrac{\sqrt3}{3}$\par
\hspace*{2.6\zw}（2）鋭角では $\sin\ang{60}=\dfrac{\sqrt3}{2}$\par
\hspace*{4.6\zw}$\sin(\ang{180}-\ang{60})=\sin\ang{60}$ なので、\ang{120}も\par
\kotae\mbox{（1）ア\enspace $\dfrac{\sqrt3}{2}$\quad イ\enspace 1\quad ウ\enspace $\dfrac{\sqrt2}{2}$\quad エ\enspace $-\dfrac{\sqrt2}{2}$\quad オ\enspace $-\dfrac{\sqrt3}{3}$}\par
\hspace*{2.6\zw}\mbox{（2）$\theta=\ang{60}$、\ang{120}}
\end{reidai}

\begin{jibun}
（1）次の値を求めましょう。\quad ア\enspace $\cos\ang{30}$\quad イ\enspace $\tan\ang{60}$\quad ウ\enspace $\sin\ang{120}$\quad エ\enspace $\cos\ang{120}$\quad オ\enspace $\tan\ang{135}$\par
（2）$\ang{0}\leqq\theta\leqq\ang{180}$ で、$\sin\theta=\dfrac{1}{2}$ になる $\theta$ を求めましょう。
\bkjans{（1）\nobreak \mbox{ア $\dfrac{\sqrt3}{2}$}、\mbox{イ $\sqrt3$}、\mbox{ウ $\dfrac{\sqrt3}{2}$}、\mbox{エ $-\dfrac{1}{2}$}、\mbox{オ $-1$}。（2）\nobreak $\theta=\ang{30}$、\ang{150}。}
\end{jibun}
\textbf{ポイント}\quad 鈍角の三角比は、$\ang{180}-\theta$ の鋭角に直して、$\sin$ はそのまま、$\cos$ と $\tan$ には $-$ をつけます。半円の図をかいて、点が $y$ 軸の右にあるか左にあるかを見れば、符号をまちがえません。

\begin{reidai}[斜面の重力を2つに分ける]
重力の大きさは $mg$（$g=\qty{9.8}{m/s^2}$）です。まさつは考えません。\par
（1）\ang{30}の斜面に、質量\qty{10}{kg}の荷物があります。斜面にそった力と、斜面へ押しつける力を求めましょう。\par
（2）傾き $\dfrac{1}{12}$ のスロープに、質量\qty{20}{kg}の台車があります。台車を止めておくのに、斜面にそって何Nの力が必要ですか。\par
\kangaekata 斜面にそった力は $mg\sin\theta$、押しつける力は $mg\cos\theta$（⑥）。（2）は、まず $\tan^{-1}$ で角を求めます（⑦）。\par
\shiki（1）$mg=10\times9.8=98$\par
\hspace*{4.6\zw}そった力 $=98\times\sin\ang{30}=98\times0.5=49$\par
\hspace*{4.6\zw}押しつける力 $=98\times\cos\ang{30}=98\times0.866=84.9$\par
\hspace*{2.6\zw}（2）$\theta=\tan^{-1}\dfrac{1}{12}\approx\ang{4.76}$、$mg=20\times9.8=196$\par
\hspace*{4.6\zw}$196\times\sin\ang{4.76}\approx196\times0.0830=16.3$\par
\kotae （1）斜面にそって\qty{49}{N}、押しつける力約\qty{84.9}{N}\quad（2）約\qty{16.3}{N}
\end{reidai}
（1）で、$49^2+84.9^2=2401+7208=9609$、$98^2=9604$ で、ほぼ等しくなります（84.9は四捨五入した値なので、少しずれる）。分けた2つの力は、長方形の2辺なので、三平方の定理が成り立つのです。$\sin^2\theta+\cos^2\theta=1$ の両辺に $(mg)^2$ をかけた形です。

\begin{jibun}
（1）\ang{45}の斜面に、質量\qty{4}{kg}の物があります。斜面にそった力と、斜面へ押しつける力を求めましょう（小数第1位まで）。\par
（2）傾き $\dfrac{1}{15}$ のスロープに、質量\qty{30}{kg}の台車があります。止めておくのに必要な、斜面にそった力を求めましょう（小数第1位まで）。
\bkjans{（1）\nobreak どちらも $39.2\times0.7071\approx\qty{27.7}{N}$。（2）\nobreak $\theta\approx\ang{3.81}$、$294\sin\ang{3.81}\approx\qty{19.6}{N}$。}
\end{jibun}
\textbf{ポイント}\quad 分けた2つの力は、どちらも重力 $mg$ より大きくなりません（$\sin\theta$ も $\cos\theta$ も1以下だから）。角が大きくなるほど、そった力はふえ、押しつける力はへります。\ang{45}で、ちょうど同じ大きさになります。

\begin{reidai}[力から斜面の角を求める]
質量\qty{5}{kg}の荷物を斜面にのせると、斜面にそって下る力が\qty{25}{N}でした。斜面の角を小数第1位まで求めましょう。まさつは考えません。\par
\kangaekata 斜面にそった力 $mg\sin\theta$ が25です。この式から $\sin\theta$ を求め、$\sin^{-1}$ で角にします。\par
\shiki $mg=5\times9.8=49$\par
\hspace*{2.6\zw}$49\times\sin\theta=25$ より\quad $\sin\theta=\dfrac{25}{49}=0.5102\cdots$\par
\hspace*{2.6\zw}$\theta=\sin^{-1}0.5102\cdots=30.68\cdots$\par
\kotae 約\ang{30.7}
\end{reidai}
確かめ：$\sin\ang{30}=0.5$ なので、\ang{30}なら $49\times0.5=24.5$。25は少しだけ大きいので、角も\ang{30}より少し大きい。合っています。

\begin{jibun}
質量\qty{8}{kg}の物を斜面にのせると、斜面にそって下る力が\qty{20}{N}でした。斜面の角を小数第1位まで求めましょう。
\bkjans{$\sin\theta=\dfrac{20}{78.4}\approx0.2551$ より約\ang{14.8}。}
\end{jibun}
\textbf{ポイント}\quad 力から角を求めるときは、まず「そった力÷重力」で $\sin\theta$ を作ります。求めた角で $mg\sin\theta$ を計算し直して、はじめの力にもどるか確かめましょう。

%% ---------------- まちがえやすいところ ----------------
\bksection{tsumazuki}{まちがえやすいところをチェック}

\begin{tsumazuki}{いつも「たて÷ななめ」を $\sin$ にする}
\noindent\begin{minipage}[c]{0.66\linewidth}
\bkng{右の図で $\sin\theta=\dfrac{\text{たて}}{\text{ななめ}}=\dfrac{4}{5}$}
\bkok{$\theta$ の向かいの辺（対辺）は横の3。$\sin\theta=\dfrac{3}{5}$}
$\sin$ は「たて」ではなく「$\theta$ の対辺」です。$\theta$ が上にあると、対辺は横になります。いつも $\theta$ の所に目を置いて、向かいの辺をさがしましょう。
\end{minipage}\hfill
\begin{minipage}[c]{0.3\linewidth}\centering
\begin{tikzpicture}[x=1mm, y=1mm]
  \coordinate (A) at (0,0); \coordinate (B) at (18,0); \coordinate (C) at (0,24);
  \draw[bkInk, line width=0.9pt] (A) -- (B) -- (C) -- cycle;
  \bkivgright{(A)}{0}{2.2}
  \draw[bkInk, line width=0.5pt] ($(C)+(270:6)$) arc (270:{270+36.87}:6);
  \node[font=\sffamily\footnotesize] at ($(C)+(288:8.8)$) {$\theta$};
  \node[font=\sffamily\footnotesize, anchor=north] at (9,-1.4) {3};
  \node[font=\sffamily\footnotesize, anchor=east] at (-1.4,12) {4};
  \node[font=\sffamily\footnotesize, anchor=south west] at (9.5,12.5) {5};
\end{tikzpicture}
\end{minipage}
\end{tsumazuki}

\begin{tsumazuki}{電卓が RAD のまま}
\bkng{$\sin30$ を計算したら $-0.988$。だから $\sin\ang{30}=-0.988$}
\bkok{DEG に切りかえて $\sin30=0.5$}
鋭角の三角比は、0より大きいはずです。$-0.988$ は「30ラジアン」のサインです。$\sin30$ で0.5が出るか、はじめに確かめる習慣をつけましょう。
\end{tsumazuki}

\begin{tsumazuki}{$\sin^{-1}$ を $\dfrac{1}{\sin}$ と読む}
\bkng{$\sin^{-1}0.5=\dfrac{1}{\sin0.5}$}
\bkok{$\sin^{-1}0.5$ は「$\sin$ が0.5になる角」で、\ang{30}}
$-1$ は「反対向きの計算」のしるしです。いっぽう $\sin^2\theta$ の2は、$(\sin\theta)^2$ という2乗です。右上の数の意味が、$-1$ と2でちがうので気をつけましょう。
\end{tsumazuki}

\begin{tsumazuki}{斜面にそった力を $mg\cos\theta$ にする}
\bkng{斜面にそった力は $mg\cos\theta$}
\bkok{斜面にそった力は $mg\sin\theta$}
$\theta=\ang{0}$（平ら）で確かめます。平らな床の上の物は動かないので、そった力は0のはず。$\cos\ang{0}=1$ では $mg$ になって合いません。$\sin\ang{0}=0$ なら合います。
\end{tsumazuki}

\begin{tsumazuki}{傾きと角度を比例させる}
\bkng{傾き1が\ang{45}なので、傾き2は\ang{90}}
\bkok{傾き2は $\tan^{-1}2\approx\ang{63.4}$}
傾きと角度は比例しません。傾き $\tan\theta$ は、\ang{90}に近づくと急に大きくなります。
\end{tsumazuki}

%% ---------------- 物理・工学 ----------------
\begin{butsuri}
\par\addvspace{3pt}\textbf{車いすのスロープ}\par\noindent
\begin{minipage}[c]{0.62\linewidth}
建物のスロープは、きまりで傾き $\dfrac{1}{12}$ より急にしないことになっています（屋外では、もっとゆるやかな $\dfrac{1}{15}$ がすすめられています）。人と車いすを合わせて質量\qty{80}{kg}なら、重力は $80\times9.8=\qty{784}{N}$。$\dfrac{1}{12}$ のスロープ（約\ang{4.76}）では、坂にそって $784\times\sin\ang{4.76}\approx\qty{65}{N}$ の力で押せばのぼれます。\ang{30}の坂なら\qty{392}{N}もかかり、とてものぼれません。
\end{minipage}\hfill
\begin{minipage}[c]{0.34\linewidth}\centering
\begin{tikzpicture}[x=0.8mm, y=0.8mm]
  \fill[bkGray!25] (0,0) -- (44,0) -- (44,{44/12}) -- cycle;
  \draw[bkInk!80, line width=0.7pt] (0,0) -- (44,0) -- (44,{44/12}) -- (52,{44/12});
  \begin{scope}[shift={(22,{22/12})}, rotate={atan(1/12)}]
    \draw[bkBlue!80!black, line width=0.8pt] (0,3.2) circle (3.2);
    \fill[bkBlue!80!black] (0,3.2) circle (0.6);
    \draw[bkBlue!80!black, line width=0.8pt] (-2,3.5) -- (-2,11) -- (-5,13);
    \draw[bkBlue!80!black, line width=0.8pt] (-2,7) -- (3,7) -- (3,3.5);
    \fill[bkSkin] (0.5,16.5) circle (2.2);
    \fill[bkGreen!70!black, rounded corners=1] (-1.6,7.5) rectangle (2.2,14);
  \end{scope}
  \node[font=\sffamily\footnotesize, anchor=north] at (22,-1.4) {12};
  \node[font=\sffamily\footnotesize, anchor=west] at (45,1.8) {1};
\end{tikzpicture}
\end{minipage}

\end{butsuri}

\begin{butsuri}
\par\addvspace{3pt}\textbf{測量：行けない所の高さをはかる}\par\noindent
\begin{minipage}[c]{0.62\linewidth}
土地の高さや建物の高さをはかる人は、トランシットという道具で仰角をはかり、\bkreidairef{4}と同じように $\tan$ で高さを計算します。スマートフォンの「高さをはかるアプリ」も、画面のかたむき（加速度センサーではかる角度）と距離から、同じ計算をしています。
\end{minipage}\hfill
\begin{minipage}[c]{0.34\linewidth}\centering
\begin{tikzpicture}[x=0.88mm, y=0.88mm]
  \fill[bkGray!30] (-2,-1) rectangle (44,0);
  \draw[bkInk, line width=0.6pt] (2,0) -- (4,9) -- (6,0) (4,0) -- (4,9);
  \fill[bkNavy] (2,9) rectangle (7,12.5);
  \fill[bkBlue!25] (32,0) rectangle (42,26);
  \draw[bkBlue!70!black, line width=0.6pt] (32,0) rectangle (42,26);
  \foreach \yy in {4,10,16,22} {\fill[white] (34,\yy) rectangle (36.5,\yy+2.5); \fill[white] (37.5,\yy) rectangle (40,\yy+2.5);}
  \draw[bkOrange, line width=0.9pt] (6,11) -- (32,26);
  \draw[bkInk, line width=0.5pt, dashed] (6,11) -- (32,11);
  \draw[bkInk, line width=0.5pt] (14,11) arc (0:30:8);
\end{tikzpicture}
\end{minipage}

\end{butsuri}

\begin{butsuri}
\par\addvspace{3pt}\textbf{すべり出す角で、まさつの大きさをはかる}\par\noindent
\begin{minipage}[c]{0.68\linewidth}
下じきに消しゴムをのせ、下じきを少しずつかたむけます。ある角 $\theta$ をこえると、消しゴムはすべり出します。まさつの力は、押しつける力 $mg\cos\theta$ の、最大で $\mu$（ミュー）倍までしか出せません。$\mu$ はすべりにくさを表す数（静止まさつ係数）です。すべり出す角では、そった力とまさつの最大がちょうど同じなので、
\bkivgline{$mg\sin\theta=\mu\,mg\cos\theta$\quad より\quad $\mu=\dfrac{\sin\theta}{\cos\theta}=\tan\theta$}
\ang{35}ですべり出したら、$\mu=\tan\ang{35}\approx0.70$ です。質量 $m$ が消えるので、重い消しゴムでも同じ角ですべり出します。
\end{minipage}\hfill
\begin{minipage}[c]{0.28\linewidth}\centering
\begin{tikzpicture}[x=0.78mm, y=0.78mm]
  \def\t{35}
  \draw[bkInk!80, line width=1.4pt] (0,0) -- ({44*cos(\t)},{44*sin(\t)});
  \draw[bkGray, line width=0.8pt] (-2,0) -- (40,0);
  \bkivgbox{({24*cos(\t)},{24*sin(\t)})}{\t}{bkBlue}
  \draw[bkInk, line width=0.5pt] (9,0) arc (0:\t:9);
  \node[font=\sffamily\footnotesize] at ({\t/2}:12) {$\theta$};
  \draw[bk arrow, bkRed, line width=1pt] ({24*cos(\t)+3*cos(\t+180)},{24*sin(\t)+3*sin(\t+180)}) ++({\t+90}:-1.2) -- ++({\t+180}:10);
\end{tikzpicture}
\end{minipage}
\end{butsuri}


\begin{column}{身のまわりの坂を、角度でくらべよう}
\begin{bkcenter}
\begin{bktable}{colspec={Q[l,m,wd=56mm] Q[c,m,wd=26mm] Q[c,m,wd=22mm]}, rowsep=0pt, row{1}={font=\sffamily\footnotesize}, cells={font=\small}}
 & 傾き（高さ÷横） & 角度 \\
車いすのスロープ（$\frac{1}{12}$） & 0.083 & 約\ang{4.8} \\
こう配10\%の道路 & 0.1 & 約\ang{5.7} \\
すべり台（\qty{120}{cm}÷\qty{240}{cm}） & 0.5 & 約\ang{26.6} \\
階段（\qty{16}{cm}÷\qty{30}{cm}） & 0.533 & 約\ang{28.1} \\
エスカレーター & 0.577 & \ang{30} \\
\end{bktable}
\end{bkcenter}
傾き（3-5）を $\tan^{-1}$ で角度に直しました。エスカレーター・階段・すべり台は、どれも\ang{30}近くです。\ang{30}の坂では重力の半分が坂にそってかかりますが、\ang{5}のスロープなら約12分の1ですみます（⑥）。
\end{column}

\begin{column}{自転車で坂をのぼると、どれくらい重い？}
人と自転車を合わせて質量\qty{70}{kg}とすると、重力は $70\times9.8=\qty{686}{N}$ です。こう配ごとに、坂にそって下る力 $686\sin\theta$ を計算しました（角は $\tan^{-1}$ で求めた）。\par\vspace{3pt}
\noindent\begin{minipage}[c]{0.5\linewidth}\centering
\begin{bktable}{colspec={Q[c,m,wd=12mm] Q[c,m,wd=16mm] Q[c,m,wd=22mm]}, colsep=2pt, rowsep=0pt, row{1}={font=\sffamily\footnotesize}, cells={font=\small}}
こう配 & 角度 & 下る力 \\
5\% & \ang{2.86} & \qty{34.3}{N} \\
10\% & \ang{5.71} & \qty{68.3}{N} \\
20\% & \ang{11.31} & \qty{134.5}{N} \\
\end{bktable}
\end{minipage}\hfill
\begin{minipage}[c]{0.47\linewidth}
こう配10\%の坂では、平らな道より約\qty{68}{N}（約\qty{7}{kg}の物を持ち上げる力）多くペダルをふみます。ギアを軽くすると、進むのはおそくなるかわりに、ふむ力が小さくてすみます。
\end{minipage}
\end{column}

\begin{column}{表計算で、三角比の表を作ろう}
\raggedright A列に角度（度）を入れ、B2に \texttt{=SIN(RADIANS(A2))}、C2に \texttt{=COS(RADIANS(A2))}、D2に \texttt{=TAN(RADIANS(A2))} と入れて、下の行へコピーします。度をラジアンに直す \texttt{RADIANS} をわすれずに（4-2）。\par\vspace{3pt}
\noindent\begin{minipage}[c]{0.6\linewidth}\centering
\bksheet{4}{14}{%
1 & 角度 & sin & cos & tan \\
2 & 0 & 0.0000 & 1.0000 & 0.0000 \\
3 & 30 & 0.5000 & 0.8660 & 0.5774 \\
4 & 45 & 0.7071 & 0.7071 & 1.0000 \\
5 & 60 & 0.8660 & 0.5000 & 1.7321 \\}
\end{minipage}\hfill
\begin{minipage}[c]{0.37\linewidth}
$\sin\ang{30}=\cos\ang{60}$（⑤）や三角定規の値が確かめられます。角度を1°ずつにすれば、「三角比の表」がすぐにできます。
\end{minipage}
\end{column}

%% ---------------- 章末問題 ----------------
\bkendproblems
\bkproblemgroup{基本}{2}

\bkproblem 直角三角形ABCで、Cが直角、$\mathrm{BC}=8$、$\mathrm{AC}=15$、$\mathrm{AB}=17$ です。$\sin A$、$\cos A$、$\tan A$ と、$\sin B$、$\cos B$、$\tan B$ を求めましょう。\mbox{（\bkreidairef{1}）}
\bkpans{$\sin A=\dfrac{8}{17}$、$\cos A=\dfrac{15}{17}$、$\tan A=\dfrac{8}{15}$、$\sin B=\dfrac{15}{17}$、$\cos B=\dfrac{8}{17}$、$\tan B=\dfrac{15}{8}$}

\bkproblem 次の値を求めましょう。\quad（1）$\sin\ang{45}$\quad（2）$\cos\ang{60}$\quad（3）$\tan\ang{30}$\quad（4）$\sin\ang{150}$\quad（5）$\cos\ang{135}$\mbox{（\bkreidairef{7}）}
\bkpans{（1）\nobreak $\dfrac{\sqrt2}{2}$。（2）\nobreak $\dfrac{1}{2}$。（3）\nobreak $\dfrac{\sqrt3}{3}$。（4）\nobreak $\dfrac{1}{2}$。（5）\nobreak $-\dfrac{\sqrt2}{2}$。}

\bkproblem 長さ\qty{300}{m}のまっすぐな坂道が、水平と\ang{6}の角をなしています。坂をのぼりきると何m高くなりますか。また、横に何m進みますか。小数第1位まで求めましょう。\mbox{（\bkreidairef{2}）}
\bkpans{高さ $300\sin\ang{6}\approx\qty{31.4}{m}$、横 $300\cos\ang{6}\approx\qty{298.4}{m}$}

\bkproblem 極座標が次の点の $(x,\ y)$ を求めましょう。（1）は $\sqrt{\phantom{2}}$ を使い、（2）は電卓で小数第1位まで求めます。\quad（1）$(6,\ \ang{45})$\quad（2）$(20,\ \ang{70})$\mbox{（\bkreidairef{3}）}
\bkpans{（1）\nobreak $x=6\cos\ang{45}=3\sqrt2$、$y=6\sin\ang{45}=3\sqrt2$ で、$(3\sqrt2,\ 3\sqrt2)$（約 $(4.24,\ 4.24)$）。（2）\nobreak $x=20\cos\ang{70}\approx6.8$、$y=20\sin\ang{70}\approx18.8$ で、約 $(6.8,\ 18.8)$。}

\bkproblem 塔から\qty{25}{m}はなれた所で、塔のてっぺんの仰角をはかると\ang{32}でした。目の高さを\qty{1.5}{m}として、塔の高さを小数第1位まで求めましょう。\mbox{（\bkreidairef{4}）}
\bkpans{$25\times\tan\ang{32}\approx25\times0.6249\approx15.62$。$15.62+1.5=17.12$ で約\qty{17.1}{m}}

\bkproblem 次の角度を、小数第1位まで求めましょう。\\
（1）傾き0.2の坂の角度\qquad（2）長さ\qty{4}{m}のはしごの上のはしが、高さ\qty{3.5}{m}にあるときの、はしごと地面の角度\mbox{（\bkreidairef{5}）}
\bkpans{（1）\nobreak $\tan^{-1}0.2\approx\ang{11.3}$。（2）\nobreak $\sin\theta=\dfrac{3.5}{4}=0.875$、$\sin^{-1}0.875\approx\ang{61.0}$}

\bkproblem $\theta$ が鋭角で、$\sin\theta=\dfrac{8}{17}$ のとき、$\cos\theta$ と $\tan\theta$ を求めましょう。\mbox{（\bkreidairef{6}）}
\bkpans{$\cos^2\theta=1-\dfrac{64}{289}=\dfrac{225}{289}$、$\cos\theta=\dfrac{15}{17}$。$\tan\theta=\dfrac{8}{15}$}

\bkproblem \ang{20}の斜面に、質量\qty{40}{kg}の荷物があります。斜面にそった力と、斜面へ押しつける力を、小数第1位まで求めましょう（$g=\qty{9.8}{m/s^2}$、まさつは考えない）。\mbox{（\bkreidairef{8}）}
\bkpans{そった力 $392\sin\ang{20}\approx\qty{134.1}{N}$、押しつける力 $392\cos\ang{20}\approx\qty{368.4}{N}$}

\bkproblemgroup{チャレンジ}{3}

\bkproblem 高さ\qty{0.75}{m}の入り口へ、傾き $\dfrac{1}{12}$ のスロープを作ります。（1）\nobreak スロープの横の長さは何mですか。（2）\nobreak スロープの角度を小数第2位まで求めましょう。（3）\nobreak スロープの板の長さ（斜辺）は何mですか。小数第2位まで求めましょう。\mbox{（\bkreidairef{5}・\bkreidairef{2}）}
\bkpans{（1）\nobreak $0.75\times12=\qty{9}{m}$。（2）\nobreak $\tan^{-1}\dfrac{1}{12}\approx\ang{4.76}$。（3）\nobreak $\dfrac{0.75}{\sin\ang{4.76}}\approx\dfrac{0.75}{0.0830}\approx\qty{9.03}{m}$（三平方の定理で $\sqrt{9^2+0.75^2}\approx9.03$ とも求められる）}

\bkproblem $\ang{0}\leqq\theta\leqq\ang{180}$ で、$\sin\theta=\dfrac{\sqrt2}{2}$ になる $\theta$ を求めましょう。また、そのそれぞれについて $\cos\theta$ の値を求めましょう。\mbox{（\bkreidairef{7}）}
\bkpans{$\theta=\ang{45}$ のとき $\cos\theta=\dfrac{\sqrt2}{2}$。$\theta=\ang{135}$ のとき $\cos\theta=-\dfrac{\sqrt2}{2}$。}

\bkproblem 質量\qty{12}{kg}の荷物を斜面にのせると、斜面にそって下る力が\qty{50}{N}でした（まさつは考えない）。（1）\nobreak 斜面の角を小数第1位まで求めましょう。（2）\nobreak 斜面へ押しつける力を小数第1位まで求めましょう。\mbox{（\bkreidairef{9}・\bkreidairef{8}）}
\bkpans{（1）\nobreak $mg=117.6$、$\sin\theta=\dfrac{50}{117.6}\approx0.4252$、$\theta\approx\ang{25.2}$。（2）\nobreak $117.6\times\cos\ang{25.2}\approx\qty{106.4}{N}$}

\begin{furikaeri}
\begin{checklist}
\item 直角三角形で、どの角から見るかに気をつけて、$\sin$・$\cos$・$\tan$ を求められる。\bkfurref{bk:4-6:teigi}
\item 三角比と電卓で、辺の長さや角度を求められる。坂の傾きを角度に直せる。\bkfurref{bk:4-6:atai}
\item 斜面の重力を、$mg\sin\theta$ と $mg\cos\theta$ に分けられる。\bkfurref{bk:4-6:bunkai}
\end{checklist}
\end{furikaeri}

\bkfukidashi{sora}{三角比があれば、行けない木の高さも、坂の上の力も計算できるね。次の章では、直角のない三角形も計算するよ。}
